\documentclass[UTF8,12pt,a4paper,fontset=fandol]{ctexbook}

% 支持从项目根目录或力学目录编译；沿用项目共享样式。
\makeatletter
\def\input@path{{../}}
\makeatother
\usepackage{researchnotes}

\title{本科力学}
\author{}
\date{}

\begin{document}
\frontmatter
\maketitle
\tableofcontents

\chapter{编写说明与学习路线}
\label{chap:mech-preface}

\section{课程定位与参考体系}
\keyterm{力学（mechanics）}研究物体的运动、相互作用以及运动变化的规律。本书面向理工科本科一年级至二年级学生，以高中物理、一元微积分和初步矢量知识为起点，建立能够独立使用的普通力学体系。核心任务是从实际情境中选择研究对象、提出可检验的近似、建立运动方程，并用守恒律、极限情形和量纲核查结论。分析力学只作为后续课程的入口，不能替代对受力、约束和物理过程的理解。

课程范围主要参照 MIT 8.01SC 的牛顿力学与守恒律主线；数学基础较好的读者可把 MIT 8.012 作为提高要求的参照，其官方课程说明将它定位为一年级力学的进阶版本，并列出 Kleppner 与 Kolenkow 的 \textit{An Introduction to Mechanics} 为教材\cite{mit801,mit8012}。MIT 8.01 在线教材的目录还提供了弹性、流体和非惯性系等后续阅读入口\cite{dourmashkin}。这些资料用于课程层次与主题覆盖的选择，正文采用独立组织的中文论述、推导和训练题，不是任何一本教材的逐章翻译，也不代表相关学校的正式教材。

中文教学组织参考清华大学陈曦的公开教学说明以及北京大学物理学院的课程体系介绍。前者强调解释假设和推理细节，同时说明其力学讲义预期读者已经具有一定基础；后者把普通物理与理论物理分层设置，并将牛顿力学和狭义相对论列入力学课程内容\cite{tsinghua-chen,pku-curriculum}。据此，本书补足初学者所需的数学和建模铺垫，将狭义相对论作为牛顿时空观的延伸，将拉格朗日方法单列为选学内容。引用资料的可核验链接列于书末，公开课程年份与本书检索日期分别标明。

\section{知识依赖与阅读顺序}
第~\ref{chap:mech-tools}--\ref{chap:mech-newton} 章建立量纲、运动学和动力学基础；第~\ref{chap:mech-energy}--\ref{chap:mech-rigid} 章从功、动量和角动量发展质点系与刚体理论；第~\ref{chap:mech-gravity}、\ref{chap:mech-frames} 章分别处理引力轨道和非惯性参考系。第~\ref{chap:mech-oscillation}、\ref{chap:mech-waves} 章由振动过渡到波，第~\ref{chap:mech-continuum} 章介绍弹性与理想流体的基础模型。第~\ref{chap:mech-relativity} 章讨论狭义相对论，第~\ref{chap:mech-lagrange} 章提供分析力学入门，第~\ref{chap:mech-computation} 章说明数值计算与实验检验的方法。

首次学习宜循序完成至振动一章，其中引力轨道的微分方程推导、惯性张量与陀螺进动可以第二遍再读。波和连续介质需要偏导数的初步知识，分析力学需要多元微分与分部积分；相关工具在使用处解释。每章末设三道自编习题，依次侧重概念、计算或推导及综合应用，附录~\ref{app:mech-solutions} 给出逐题解答。例题用于示范建模过程，习题用于检验能否迁移方法，两者都应先明确条件再计算。

\section{符号、假设与论证约定}
空间矢量用粗体小写或大写字母，如位置 $\mathbf r$、速度 $\mathbf v$、力 $\mathbf F$；对应的不带粗体符号通常表示大小，但坐标分量另加下标。时间记为 $t$，上点表示对时间求导。采用右手直角坐标系和国际单位制（International System of Units, SI），角度在微分公式中以弧度计。$m$ 表示单个物体的质量，$M$ 常表示系统总质量，$g>0$ 表示局部重力加速度的大小。局部符号的含义以首次定义为准；例如 $T$ 在不同章节中可表示周期、拉力或动能，不在同一推导内混用。

除特别说明，牛顿力学部分研究质量恒定的物质系统，时间连续，位置至少分段二阶可微，参考系近似惯性。碰撞采用短时冲量模型时允许速度在理想极限中跳变；摩擦方向改变和接触脱离需要分段建模。定理是在给定物理模型内的数学结论；牛顿定律、引力定律和材料本构关系则属于需要实验支持的物理规律，不能仅由数学证明。数值结果是离散模型的近似，不能仅凭图像宣布守恒、稳定或收敛。

\mainmatter

\chapter{物理量、测量与数学工具}
\label{chap:mech-tools}

\section{从物体到力学模型}
一辆汽车经过路口时，研究到达时间可以忽略车身内部结构；研究转弯稳定性则必须考虑车身大小和质量分布。\keyterm{质点（particle）}是只保留质量与位置的物体模型，适用性由研究问题决定，不由物体绝对尺寸决定。\keyterm{刚体（rigid body）}模型要求任意两物质点间距离保持不变，适合形变相对较小且不影响目标量的情形；弹性体和流体则允许形变。

建模需要同时说明研究系统、参考系、相互作用、约束和初始条件。例如忽略空气阻力的抛体模型，不是宣称空气不存在，而是假设在所考察时间与速度范围内，阻力造成的变化低于所需精度。若初步估计阻力 $F_{\mathrm d}$ 与重力之比 $F_{\mathrm d}/(mg)$ 已接近 $1$，就不能继续用该近似解释整个过程。模型的精细程度必须与问题和证据相配。

\section{量纲与数量级}
物理量的数值与单位必须同时给出。用 $\mathsf L,\mathsf M,\mathsf T$ 分别表示长度、质量和时间的量纲，则
\[
 [v]=\mathsf L\mathsf T^{-1},\quad [a]=\mathsf L\mathsf T^{-2},\quad
 [F]=\mathsf M\mathsf L\mathsf T^{-2},\quad [E]=\mathsf M\mathsf L^2\mathsf T^{-2}.
\]
等式两端和被相加的各项必须量纲一致；指数、对数和三角函数的自变量应无量纲。量纲正确是公式成立的必要条件，但不能决定无量纲常数、正负号或全部函数形式。

\begin{example}[用量纲限制单摆周期]
假设小摆幅、无阻力单摆的周期 $T$ 仅依赖摆长 $\ell$、重力加速度 $g$ 和摆球质量 $m$，且依赖形式为 $T=C\ell^a g^b m^d$，其中 $C$ 无量纲。求指数。
\end{example}
\begin{solve}
比较质量、时间和长度的幂，得 $d=0$、$-2b=1$、$a+b=0$，故 $T=C\sqrt{\ell/g}$。这说明模型中的周期与质量无关，但 $C=2\pi$ 需要动力学推导。若摆幅 $\theta_0$ 不能忽略，量纲分析只给出 $T=\sqrt{\ell/g}\,f(\theta_0)$，因为 $\theta_0$ 无量纲。
\end{solve}

\section{矢量、投影与叉积}
\keyterm{矢量（vector）}既有大小又有方向，其分量随坐标基的变化而变化，但几何对象本身不变。在正交单位基 $\mathbf e_x,\mathbf e_y,\mathbf e_z$ 下，$\mathbf a=a_x\mathbf e_x+a_y\mathbf e_y+a_z\mathbf e_z$。点积和叉积定义为
\[
 \mathbf a\cdot\mathbf b=ab\cos\theta=\sum_i a_i b_i,\qquad
 \mathbf a\times\mathbf b=
 \begin{vmatrix}\mathbf e_x&\mathbf e_y&\mathbf e_z\\a_x&a_y&a_z\\b_x&b_y&b_z\end{vmatrix}.
\]
叉积方向由右手规则确定，其大小为 $ab\sin\theta$。点积刻画沿给定方向的投影，叉积刻画垂直于两矢量平面的有向面积；功和力矩分别采用这两种结构。不能把矢量方程中方向不同的力直接按大小相加。

若基矢固定，则矢量求导逐分量进行；若基矢随时间转动，必须同时求基矢的导数。对可微矢量，乘法法则仍成立，例如
\[
 \frac{\mathrm d}{\mathrm dt}(\mathbf a\cdot\mathbf b)=\dot{\mathbf a}\cdot\mathbf b+\mathbf a\cdot\dot{\mathbf b},\qquad
 \frac{\mathrm d}{\mathrm dt}(\mathbf a\times\mathbf b)=\dot{\mathbf a}\times\mathbf b+\mathbf a\times\dot{\mathbf b}.
\]
叉积中的顺序不可交换，因为 $\mathbf a\times\mathbf b=-\mathbf b\times\mathbf a$。

\section{微分方程、初值和近似}
运动规律通常给出加速度，而要求解的是位置。方程 $\ddot x=a$ 的一般解含两个常数，初始位置 $x(0)=x_0$ 和速度 $\dot x(0)=v_0$ 才将其确定。更一般地，若 $\ddot{\mathbf r}=\mathbf f(\mathbf r,\dot{\mathbf r},t)$ 的右端对 $t$ 连续，对位置和速度局部 Lipschitz 连续，则每组初值在足够小的时间区间内有唯一解。这一常微分方程结论解释了初值的作用，但不保证解永远存在，也不适用于未经处理的瞬时碰撞。

近似式应带有控制参数。若 $|\varepsilon|\ll1$，则 $(1+\varepsilon)^\alpha=1+\alpha\varepsilon+O(\varepsilon^2)$，其中 $O(\varepsilon^2)$ 表示误差绝对值在足够小的邻域内被常数乘 $\varepsilon^2$ 控制。小角度近似为 $\sin\theta=\theta+O(\theta^3)$；它要求弧度下的 $|\theta|$ 小，不能仅凭角度数值写成整数就判定小量。

测量值也带有不确定度。对 $y=f(x_1,\ldots,x_n)$，当误差较小、$f$ 在测量点附近光滑，且各输入的随机误差互不相关时，一阶传播给出
\begin{equation}
 u_y^2\approx\sum_{i=1}^n\left(\frac{\partial f}{\partial x_i}\right)^2u_i^2.
 \label{eq:mech-uncertainty}
\end{equation}
这里 $u_i$ 是标准不确定度；有关联时应加入协方差项。系统偏差不会因为重复测量就自动消失，公式也不能替代对仪器、实验步骤和模型偏差的分析。

\section{习题}
\begin{enumerate}[label=\thechapter.\arabic*,ref=\thechapter.\arabic*]
\item\label{exr:mech-dimension} 设线性阻力为 $\mathbf F=-b\mathbf v$，平方阻力为 $\mathbf F=-c_{\mathrm d}|\mathbf v|\mathbf v$。求 $b,c_{\mathrm d}$ 的量纲，并解释 $m/b$ 的物理量纲。
\item\label{exr:mech-cross} 设 $\mathbf r=(2,0,0)\,\mathrm m$，$\mathbf F=(0,3,0)\,\mathrm N$。求 $\mathbf r\cdot\mathbf F$ 与 $\mathbf r\times\mathbf F$，说明前者为什么不能单凭量纲就解释为实际做功。
\item\label{exr:mech-error} 用 $g=4\pi^2\ell/T^2$ 测量重力加速度。若 $\ell,T$ 的标准不确定度相互独立，推导 $u_g/g$，并说明摆幅过大属于哪一类误差来源。
\end{enumerate}

\chapter{质点运动学}
\label{chap:mech-kinematics}

\section{参考系、轨迹与速度}
\keyterm{参考系（reference frame）}包括用于描述位置的空间基准和记录时间的时钟约定；坐标系是同一参考系中标记位置的数学方法。换用极坐标并不自动意味着换了参考系。给定参考系后，位置矢量 $\mathbf r(t)$ 描述运动，轨迹是其在空间中经过的点集，同一条轨迹可以对应不同的时间规律。

\begin{definition}[位移、速度与加速度]
在时间区间 $[t_1,t_2]$ 内，位移为 $\Delta\mathbf r=\mathbf r(t_2)-\mathbf r(t_1)$。若 $\mathbf r$ 可微，瞬时速度为 $\mathbf v=\mathrm d\mathbf r/\mathrm dt$，速率为 $v=|\mathbf v|$；若再可微，加速度为 $\mathbf a=\mathrm d\mathbf v/\mathrm dt$。
\end{definition}
路程为 $s=\int_{t_1}^{t_2}|\mathbf v(t)|\,\mathrm dt$，通常大于位移的大小。平均速率 $s/(t_2-t_1)$ 与平均速度 $\Delta\mathbf r/(t_2-t_1)$ 是不同对象。速度方向变化也产生加速度，因此匀速率运动未必是匀速直线运动。

对已知连续加速度，积分形式为
\begin{equation}
 \mathbf v(t)=\mathbf v_0+\int_0^t\mathbf a(\tau)\,\mathrm d\tau,\qquad
 \mathbf r(t)=\mathbf r_0+\int_0^t\mathbf v(\tau)\,\mathrm d\tau.
 \label{eq:mech-kinematic-integrals}
\end{equation}
常加速度时才可简化为 $\mathbf r=\mathbf r_0+\mathbf v_0t+\mathbf a t^2/2$。一维公式 $v^2-v_0^2=2a(x-x_0)$ 中的 $a$ 必须恒定；对一般位置依赖加速度，可在速度不为零的单调运动区间使用 $a=v\,\mathrm dv/\mathrm dx$。

\section{抛体运动与相遇条件}
取水平 $x$ 轴和竖直向上 $y$ 轴，在均匀重力场中忽略阻力。由 $\mathbf a=(0,-g)$ 得
\[
 x=x_0+v_{0x}t,\qquad y=y_0+v_{0y}t-\tfrac12gt^2.
\]
当 $v_{0x}\ne0$ 时消去 $t$ 得抛物线，但物理轨迹仅包含满足时间范围和接触条件的部分。落地、碰墙等事件由位置条件求得；不能先使用同高度射程公式再修改落点高度。

\begin{example}[从高处水平抛出]
从离水平地面高度 $h>0$ 处，以水平速度 $u>0$ 抛出质点，忽略阻力，求落地时间、水平距离和落地前速度。
\end{example}
\begin{solve}
取抛出位置的竖直投影为原点，则 $x=ut$、$y=h-gt^2/2$。首次到达 $y=0$ 的正时间为 $t_*=\sqrt{2h/g}$，所以 $R=u\sqrt{2h/g}$，$\mathbf v(t_*)=(u,-\sqrt{2gh})$，速率为 $\sqrt{u^2+2gh}$。这里的时间与水平初速度无关，来自水平、竖直分量在该模型中的独立演化。
\end{solve}

\section{曲线运动与极坐标}
对正则轨迹，在速率非零的区间定义单位切向量 $\mathbf e_{\mathrm t}=\mathbf v/v$。若曲率为 $\kappa=1/\rho$，主法向单位矢量为 $\mathbf e_{\mathrm n}$，则 $\mathrm d\mathbf e_{\mathrm t}/\mathrm ds=\mathbf e_{\mathrm n}/\rho$，从而
\begin{equation}
 \mathbf a=\dot v\,\mathbf e_{\mathrm t}+\frac{v^2}{\rho}\mathbf e_{\mathrm n}.
 \label{eq:mech-tangent-normal}
\end{equation}
切向加速度改变速率，法向加速度改变方向。直线运动可将曲率取为零；速度为零时，切向方向需要由具体轨迹或极限另行确定。

平面极坐标写作 $\mathbf r=r\mathbf e_r$，其中 $\mathbf e_r=(\cos\theta,\sin\theta)$、$\mathbf e_\theta=(-\sin\theta,\cos\theta)$。由于 $\dot{\mathbf e}_r=\dot\theta\mathbf e_\theta$、$\dot{\mathbf e}_\theta=-\dot\theta\mathbf e_r$，连续求导得到
\begin{align}
 \mathbf v&=\dot r\mathbf e_r+r\dot\theta\mathbf e_\theta,\label{eq:mech-polar-v}\\
 \mathbf a&=(\ddot r-r\dot\theta^2)\mathbf e_r+(r\ddot\theta+2\dot r\dot\theta)\mathbf e_\theta.
 \label{eq:mech-polar-a}
\end{align}
公式适用于 $r>0$ 的可微运动。径向分量一般不是 $\ddot r$；只有角速度为零等特殊情形才可这样写。圆周运动有 $r=R$，于是 $a_r=-R\dot\theta^2$、$a_\theta=R\ddot\theta$。

\begin{example}[匀速转动杆上的滑珠运动学]
一根直杆绕固定端在水平面内以恒定角速度 $\omega$ 转动。若珠沿杆以恒定相对速率 $u$ 向外运动，且 $r=r_0+ut>0$，求其速度与加速度。
\end{example}
\begin{solve}
由式~\eqref{eq:mech-polar-v} 和式~\eqref{eq:mech-polar-a}，得 $\mathbf v=u\mathbf e_r+\omega(r_0+ut)\mathbf e_\theta$，$\mathbf a=-\omega^2(r_0+ut)\mathbf e_r+2u\omega\mathbf e_\theta$。珠沿杆的相对速率恒定，并不意味着它在地面参考系中的速率或速度恒定。此处只规定运动，维持该运动所需的力需用动力学另求。
\end{solve}

\section{相对运动}
设参考系 $S'$ 的原点在 $S$ 中的位置为 $\mathbf R(t)$，两系坐标轴始终平行，则 $\mathbf r=\mathbf R+\mathbf r'$，故
\[
 \mathbf v=\dot{\mathbf R}+\mathbf v',\qquad \mathbf a=\ddot{\mathbf R}+\mathbf a'.
\]
当 $\dot{\mathbf R}=\mathbf V$ 为常矢量时，这就是\keyterm{伽利略变换（Galilean transformation）}的速度和加速度关系，且时间相同。它是牛顿时空观中的变换规律；相对速度接近光速时需使用第~\ref{chap:mech-relativity} 章的变换。

\section{习题}
\begin{enumerate}[label=\thechapter.\arabic*,ref=\thechapter.\arabic*]
\item\label{exr:mech-path} 质点沿直线运动，$x(t)=At^2-Bt^3$，其中 $A,B>0$。求 $0\le t\le A/B$ 内的位移和路程。
\item\label{exr:mech-circle} 质点沿半径 $R$ 的圆运动，$\theta(t)=\alpha t^2/2$，其中 $\alpha>0$。求加速度大小，并求切向与法向加速度大小相等的正时间。
\item\label{exr:mech-river} 河水以速度 $u$ 沿岸流动，船相对水的速率恒为 $v>u$，河宽 $d$。分别求最短渡河时间和到达正对岸所需时间，说明所选航向。
\end{enumerate}

\chapter{牛顿定律与约束动力学}
\label{chap:mech-newton}

\section{惯性系与运动定律}
\keyterm{惯性参考系（inertial reference frame）}是自由质点保持静止或匀速直线运动的参考系。实际判断需要将多种近似隔离的物体和相互作用模型结合检验，不能因某个受力物体恰好静止就认定所在系是惯性系。地面在许多短时、小尺度问题中可近似为惯性系，精密或长时间实验则可能需要考虑地球自转。

牛顿第一定律明确惯性运动及惯性系的角色；第二定律在惯性系中将动量变化与合外力联系起来。对于质量恒定的质点，\keyterm{动量（momentum）}为 $\mathbf p=m\mathbf v$，于是
\begin{equation}
 \frac{\mathrm d\mathbf p}{\mathrm dt}=\mathbf F_{\mathrm{net}},\qquad
 m\mathbf a=\sum_j\mathbf F_j.
 \label{eq:mech-newton}
\end{equation}
该关系是物理定律，不能仅用 $\mathbf F=m\mathbf a$ 作为力的定义后宣称已经验证定律。质量恒定也很关键，吸收或排出物质的系统必须另计动量通量。

牛顿第三定律的弱形式是两物体相互作用力等大反向，$\mathbf F_{ij}=-\mathbf F_{ji}$，其中 $\mathbf F_{ij}$ 表示 $j$ 对 $i$ 的力；沿两质点连线的强形式还能保证内力矩成对抵消。本书的引力、理想弹簧等模型满足相应条件。一般电磁相互作用若只考察粒子而不包含场，不一定满足这种瞬时质点对力模型，不能无条件推广。

\section{常见力与受力分析}
近地面重力近似为 $\mathbf F_g=-mg\mathbf e_y$。理想线性弹簧沿自身方向施力，一维形式为 $F_x=-kx$，其中 $x$ 从自然长度对应位置量起，$k>0$ 为劲度系数（spring constant）；这一关系仅在材料的线性弹性范围内适用。绳张力沿绳方向且只能拉不能推；质量可忽略、不可伸长的绳经无摩擦、转动惯量可忽略的滑轮时，各段张力相同，任一条件改变都应重新分析。

\keyterm{支持力（normal force）}是接触力的法向分量，不预设等于 $mg$。\keyterm{静摩擦（static friction）}在简化库仑模型中满足 $|f_s|\le\mu_s N$，实际大小由不发生相对滑动的约束决定；\keyterm{动摩擦（kinetic friction）}近似满足 $|f_k|=\mu_kN$，方向反对接触面间的相对滑动。摩擦不总是反对物体相对地面的运动，也不总是耗散该物体的机械能。

受力分析先隔离一个明确的物体，再逐一写出其他物体对它的作用。作用反作用力作用在不同物体上，不应在单个物体的受力图中相互抵消。“向心力”是合力指向曲率中心的分量名称，不是在重力、支持力之外多加的一种力。

\section{约束与接触条件}
\keyterm{约束（constraint）}限制允许的运动，例如质点沿固定光滑轨道运动或两端由不可伸长的绳连接。几何关系求导后给出速度和加速度关系，应与牛顿方程联立。接触力 $N\ge0$、张力 $T\ge0$ 属于额外的物理可行条件；算出负值通常意味着原先假定的接触或绳绷紧状态不能维持。

\begin{example}[阿特伍德装置]
质量 $m_2>m_1>0$ 的两物体通过理想绳和理想滑轮连接，系统从静止释放，求加速度和张力。
\end{example}
\begin{solve}
分别取 $m_2$ 向下、$m_1$ 向上为各自正方向，绳长约束使二者加速度大小同为 $a$。由 $m_2g-T=m_2a$、$T-m_1g=m_1a$，得
\[
 a=\frac{m_2-m_1}{m_1+m_2}g,\qquad T=\frac{2m_1m_2}{m_1+m_2}g.
\]
当两质量接近相等时 $a\to0$；当 $m_1\to0$ 时 $a\to g$、$T\to0$，与近似自由落体相符。非理想滑轮两侧张力一般不同，见第~\ref{chap:mech-rigid} 章。
\end{solve}

\begin{example}[斜面上的摩擦状态]
质量 $m$ 的物体置于固定倾角 $0<\alpha<\pi/2$ 的斜面上，静、动摩擦系数分别为 $\mu_s,\mu_k$。判断静止条件，并求物体向下滑动时的加速度。
\end{example}
\begin{solve}
法向无加速度，故 $N=mg\cos\alpha$。若静止，需静摩擦 $f_s=mg\sin\alpha$ 向上，因此可静止的条件是 $\tan\alpha\le\mu_s$。已知物体向下滑动时，沿斜面向下有 $a=g(\sin\alpha-\mu_k\cos\alpha)$。当该值为负，物体在下滑过程中减速；到速度为零后必须重新判断静摩擦是否足以维持静止，不能把原方程直接延伸为反向运动。
\end{solve}

\section{阻力与终端速度}
考虑从静止竖直下落的物体，取向下为正，在线性阻力模型下
\[
 m\dot v=mg-bv,\qquad v(0)=0,\qquad b>0.
\]
令时间尺度 $\tau=m/b$、终端速率 $v_\infty=mg/b$，将方程乘以积分因子 $e^{t/\tau}$，得
\begin{equation}
 v(t)=v_\infty(1-e^{-t/\tau}),\qquad
 y(t)-y(0)=v_\infty\bigl[t-\tau(1-e^{-t/\tau})\bigr].
 \label{eq:mech-drag}
\end{equation}
短时间 $t\ll\tau$ 时恢复 $v\approx gt$，长时间趋向终端速度。实际阻力规律与物体形状、流速和流体性质有关，线性与平方模型不能在任意速度范围内互换；拟合某个终端速度相同也不意味着全过程相同。

\section{习题}
\begin{enumerate}[label=\thechapter.\arabic*,ref=\thechapter.\arabic*]
\item\label{exr:mech-elevator} 电梯以向上加速度 $a$ 运动，质量 $m$ 的人保持接触地板，求地板支持力，并说明 $a=-g$ 的极限含义。
\item\label{exr:mech-pull} 在粗糙水平面上，物体受大小 $F$、与水平成向上角 $\alpha$ 的拉力而向右滑动。求加速度与保持接触的条件，动摩擦系数为 $\mu_k$。
\item\label{exr:mech-drag-initial} 在线性阻力下落模型中，将初速度改为任意向下分量 $v_0$。求 $v(t)$，并讨论 $v_0>v_\infty$ 和 $v_0<0$ 时的运动趋势。
\end{enumerate}

\chapter{功、能量与势能}
\label{chap:mech-energy}

\section{功和动能定理}
力与位移的联系提供了求速度的另一条道路。\keyterm{功（work）}是沿实际运动路径 $C$ 的线积分：
\begin{equation}
 W_{A\to B}=\int_C\mathbf F\cdot\mathrm d\mathbf r
 =\int_{t_A}^{t_B}\mathbf F\cdot\mathbf v\,\mathrm dt.
 \label{eq:mech-work}
\end{equation}
功依赖力的方向和实际作用点的位移，单位为焦耳 $\mathrm J=\mathrm{N\,m}$。\keyterm{功率（power）}为 $P=\mathrm dW/\mathrm dt=\mathbf F\cdot\mathbf v$，单位为瓦特。恒力公式 $W=Fs\cos\theta$ 是线积分的特殊情形，不能把变力的末值乘路程当作总功。

\begin{theorem}[质点动能定理]
在惯性系中，设质点质量 $m$ 恒定，运动满足牛顿方程且相关积分存在。以 $K=mv^2/2$ 定义\keyterm{动能（kinetic energy）}，则合力做功等于动能增量：
\begin{equation}
 K_B-K_A=W_{\mathrm{net},A\to B}.
 \label{eq:mech-work-energy}
\end{equation}
\end{theorem}
\begin{proof}
将式~\eqref{eq:mech-newton} 与速度点乘，得 $\mathbf F_{\mathrm{net}}\cdot\mathbf v=m\dot{\mathbf v}\cdot\mathbf v=\mathrm d(mv^2/2)/\mathrm dt$。对时间积分即得。速度可以改变方向，证明不要求直线运动。
\end{proof}

动能定理无需合力保守，含摩擦、阻力也成立。对于多个物体，总功应按每个力作用点的位移计算；即使一对内力等大反向，只要作用点位移不同，它们的总功也未必为零。

\section{保守力与势能}
\begin{definition}[保守力和势能]
在连通区域 $D$ 中，若不显含时间的力场 $\mathbf F(\mathbf r)$ 沿任意两点间分段光滑路径做功与路径无关，称其为\keyterm{保守力（conservative force）}。选参考点 $\mathbf r_0$，定义\keyterm{势能（potential energy）}
\[
 U(\mathbf r)-U(\mathbf r_0)=-\int_{\mathbf r_0}^{\mathbf r}\mathbf F\cdot\mathrm d\mathbf l.
\]
当 $U$ 可微时，$\mathbf F=-\nabla U$，其中 $\nabla U=(\partial_xU,\partial_yU,\partial_zU)$。
\end{definition}
势能只在加法常数意义下确定，可观测的功依赖势能差。若力场一阶连续可微，存在势能必有 $\nabla\times\mathbf F=0$；反过来，零旋度在单连通区域内足以保证保守性，区域有孔洞时不能省略全局检查。

匀强重力场中 $U=mgy$，弹簧势能为 $U=kx^2/2$。重力势能属于物体与引力源的相互作用；把引力源固定后写成质点位置的函数，是所选近似模型的简化。

\begin{proposition}[机械能变化]
设恒质量质点受力为 $\mathbf F=-\nabla U(\mathbf r,t)+\mathbf F_{\mathrm{nc}}$，且 $U$ 一阶可微，则\keyterm{机械能（mechanical energy）} $E=K+U$ 满足
\begin{equation}
 \frac{\mathrm dE}{\mathrm dt}=\frac{\partial U}{\partial t}+\mathbf F_{\mathrm{nc}}\cdot\mathbf v.
 \label{eq:mech-energy-rate}
\end{equation}
特别地，若 $U$ 不显含时间且非保守力的总功率为零，则机械能守恒。
\end{proposition}
\begin{proof}
由动能定理的微分形式及链式法则，$\dot K=(-\nabla U+\mathbf F_{\mathrm{nc}})\cdot\mathbf v$，$\dot U=\nabla U\cdot\mathbf v+\partial_tU$。两式相加即得。
\end{proof}

机械能不守恒并不意味着包含内能在内的总能量不守恒。摩擦常把有序运动的能量转化为内能；电机和外部操纵机构也可向所选系统输入能量。守恒与否首先取决于系统边界及模型中保留的能量形式。

\section{势能曲线、转折点与稳定性}
对于一维自治保守系统，$m\dot x^2/2+U(x)=E$，允许区域满足 $E\ge U(x)$；满足 $E=U(x)$ 的点可能是转折点。若 $U'(x_*)=0$，$x_*$ 是平衡位置。设 $U$ 在其附近二阶可微，则
\[
 U(x_*+\xi)=U(x_*)+\tfrac12U''(x_*)\xi^2+o(\xi^2).
\]
当 $U''(x_*)>0$ 时，线性化方程为 $m\ddot\xi=-U''(x_*)\xi$，对应稳定的小振动；当 $U''(x_*)<0$ 时，小扰动存在指数增长方向。若二阶导数为零，这个判据不作结论，必须考察高阶项，例如 $U=x^4$ 与 $U=-x^4$ 在原点具有不同稳定性。

\begin{example}[弹簧压缩与摩擦耗能]
质量 $m$ 的物块以速率 $v_0>0$ 在粗糙水平面上接触自然长度的弹簧，劲度系数为 $k$，动摩擦系数为 $\mu_k$。求首次停下时的最大压缩量 $x_*>0$。
\end{example}
\begin{solve}
在压缩阶段位移向前，摩擦做功 $-\mu_kmgx_*$，故 $mv_0^2/2=kx_*^2/2+\mu_kmgx_*$。取正根，得
\[
 x_* =\frac{\sqrt{(\mu_kmg)^2+kmv_0^2}-\mu_kmg}{k}.
\]
该结论只描述首次停下的位置。随后是否反弹取决于 $kx_*$ 是否超过最大静摩擦力，不能由压缩阶段的能量方程独自判断。
\end{solve}

\begin{example}[竖直圆轨道的接触条件]
质点在半径 $R$ 的固定光滑竖直圆形轨道内侧运动，从最低点以速率 $v_0$ 出发。求在轨道约束下完成整周运动的最低初速。
\end{example}
\begin{solve}
在最高点，向心方程为 $mg+N=mv_{\mathrm{top}}^2/R$，接触要求 $v_{\mathrm{top}}^2\ge gR$。最高点比最低点高 $2R$，故能量守恒给出 $v_0^2=v_{\mathrm{top}}^2+4gR\ge5gR$。以从最低点起的角度 $\theta$ 表示位置，沿轨道还有 $N/m=v_0^2/R-2g+3g\cos\theta$，它在最高点最小，因此最低初速为 $\sqrt{5gR}$。若质点被双向约束在刚性圆环上，允许轨道向内或向外施力，条件会改变。
\end{solve}

\section{习题}
\begin{enumerate}[label=\thechapter.\arabic*,ref=\thechapter.\arabic*]
\item\label{exr:mech-potential} 一维势能为 $U(x)=ax^4-bx^2$，其中 $a,b>0$。求全部平衡点并判断小扰动稳定性。
\item\label{exr:mech-power} 质量 $m$ 的物体在光滑水平面上由恒功率 $P>0$ 的驱动装置加速，初速 $v_0>0$。求 $v(t)$ 和位移，说明为什么不能直接把 $v_0=0$ 作为有限驱动力启动模型。
\item\label{exr:mech-timed-u} 给定 $U(x,t)=k(t)x^2/2$，无其他力。推导机械能变化率，并解释缓慢改变 $k(t)$ 是否足以保证机械能严格守恒。
\end{enumerate}

\chapter{动量、质心与变质量系统}
\label{chap:mech-momentum}

\section{冲量与系统动量}
\keyterm{冲量（impulse）}定义为 $\mathbf J=\int_{t_1}^{t_2}\mathbf F\,\mathrm dt$。对恒质量质点积分牛顿方程，得 $\Delta\mathbf p=\mathbf J$。冲量强调力的时间累积，而功强调力与位移的关系，两者不可以互相替代。

考虑由固定的一组 $N$ 个恒质量质点组成的\keyterm{封闭物质系统（closed material system）}，其总动量为 $\mathbf P=\sum_i m_i\mathbf v_i$。若内力成对满足 $\mathbf F_{ij}=-\mathbf F_{ji}$，求和消去内力即得
\begin{equation}
 \dot{\mathbf P}=\mathbf F_{\mathrm{ext}},\qquad
 \Delta\mathbf P=\mathbf J_{\mathrm{ext}}.
 \label{eq:mech-system-momentum}
\end{equation}
“封闭”指不交换物质，并不表示不受外力。总外力恒为零时，总动量在整个过程中守恒；只有总外冲量为零时，可以断言两个端点的动量相同，不能据此断言中间时刻恒定。一个方向的外力为零，足以保证该方向的动量分量守恒。

\section{质心与动能分解}
\begin{definition}[质心]
对总质量 $M=\sum_i m_i>0$ 的质点系，\keyterm{质心（center of mass）}为
\[
 \mathbf R=\frac1M\sum_i m_i\mathbf r_i.
\]
连续分布时改为 $\mathbf R=M^{-1}\int\mathbf r\,\mathrm dm$；若体密度为 $\rho$，则 $\mathrm dm=\rho\,\mathrm dV$。
\end{definition}
恒质量系统满足 $\mathbf P=M\dot{\mathbf R}$，由式~\eqref{eq:mech-system-momentum} 得 $M\ddot{\mathbf R}=\mathbf F_{\mathrm{ext}}$。质心可位于物体外部或中空区域，它是质量分布的几何特征，不必对应某个物质点。

令 $\mathbf v_i=\mathbf V+\mathbf v_i'$，其中 $\mathbf V=\dot{\mathbf R}$，则 $\sum_i m_i\mathbf v_i'=0$。展开平方得到\keyterm{柯尼希动能分解（König's theorem）}
\begin{equation}
 K=\sum_i\tfrac12m_i v_i^2=\tfrac12MV^2+\sum_i\tfrac12m_i(v_i')^2.
 \label{eq:mech-koenig}
\end{equation}
第一项是整体平动，第二项是相对质心的运动。即使质心静止，物体仍可转动、振动或含有内部热运动，因此总动能不一定为零。

\section{碰撞与恢复系数}
\keyterm{碰撞（collision）}是相互作用在相对较短时间内造成显著动量变化的过程。若外冲量相对碰撞冲量可忽略，可以近似使用动量守恒；重力并未消失，只是其在短碰撞时段的累积可忽略。\keyterm{弹性碰撞（elastic collision）}还保持碰撞前后的总动能，完全非弹性碰撞则使两物体碰后粘合。

对一维迎头碰撞，设碰前 $u_1>u_2$，碰后速度为 $v_1,v_2$。\keyterm{恢复系数（coefficient of restitution）}定义为 $e=(v_2-v_1)/(u_1-u_2)$。无额外内部能释放的常用被动碰撞模型取 $0\le e\le1$。联立动量守恒与恢复关系，有
\begin{align}
 v_1&=\frac{m_1-em_2}{m_1+m_2}u_1+\frac{(1+e)m_2}{m_1+m_2}u_2,\label{eq:mech-collision1}\\
 v_2&=\frac{(1+e)m_1}{m_1+m_2}u_1+\frac{m_2-em_1}{m_1+m_2}u_2.
\end{align}
定义\keyterm{约化质量（reduced mass）} $\mu=m_1m_2/(m_1+m_2)$，由质心动能分解可得
\begin{equation}
 K_{\mathrm{before}}-K_{\mathrm{after}}=\tfrac12\mu(1-e^2)(u_1-u_2)^2.
 \label{eq:mech-collision-loss}
\end{equation}
恢复系数是碰撞材料与速度等因素的经验描述，不是所有物体普遍相同的常数。多维斜碰需要同时处理法向、切向冲量与转动，不能仅套一维式。

\begin{example}[弹道摆的两个阶段]
质量 $m$ 的弹丸以水平速度 $u$ 射入质量 $M$ 的悬挂木块并嵌入，碰撞时间很短，碰撞时悬绳竖直。整体随后上升高度 $h$，忽略摆动阶段阻力，求 $u$。
\end{example}
\begin{solve}
碰撞阶段水平外冲量可忽略，故 $mu=(M+m)V$；竖直绳张力可能有冲量，但不改变此时的水平动量平衡。摆动阶段机械能守恒，$\tfrac12(M+m)V^2=(M+m)gh$，所以 $u=(M+m)\sqrt{2gh}/m$。嵌入过程是非弹性碰撞，若跨越两个阶段直接令弹丸初动能等于最终重力势能，会得到错误结果。
\end{solve}

\section{火箭方程与质量通量}
对不断排气的火箭，“当前火箭”并不是固定的一组物质，不能直接把 $\mathrm d(mv)/\mathrm dt$ 等同于外力。取某时刻的火箭及随后极短时间内排出的燃气为封闭系统。设火箭瞬时质量、速度为 $m,v$，在 $\mathrm dt$ 内质量变化 $\mathrm dm<0$，排气相对火箭向后的速度大小为 $u_e>0$。忽略二阶小量，动量平衡为
\[
 (m+\mathrm dm)(v+\mathrm dv)+(-\mathrm dm)(v-u_e)-mv=F_{\mathrm{ext}}\,\mathrm dt,
\]
因而
\begin{equation}
 m\dot v=F_{\mathrm{ext}}-u_e\dot m.
 \label{eq:mech-rocket}
\end{equation}
若 $F_{\mathrm{ext}}=0$ 且 $u_e$ 恒定，积分得 $v_f-v_i=u_e\ln(m_i/m_f)$。推力来自排出质量的动量通量，不需要空气提供反作用；实际喷管压力等影响可在更完整的推力模型中加入。

\begin{example}[静止雨滴落入小车]
小车在水平无摩擦轨道上运动，初质量 $m_0$、初速度 $v_0$。雨滴在地面系中无水平速度，以恒定质量流率 $q>0$ 落入并随车运动，求车速。
\end{example}
\begin{solve}
在 $\mathrm dt$ 内加入的小雨滴水平动量为零，因此 $(m+q\mathrm dt)(v+\mathrm dv)=mv$ 到一阶精度给出 $m\dot v=-qv$。又 $m=m_0+qt$，故 $v=m_0v_0/(m_0+qt)$。车与已经收入的雨水的动量数值在这一特殊输入条件下保持不变，但动能 $K=m_0^2v_0^2/[2(m_0+qt)]$ 下降；吸收过程使相对运动能转化为内能。
\end{solve}

\section{习题}
\begin{enumerate}[label=\thechapter.\arabic*,ref=\thechapter.\arabic*]
\item\label{exr:mech-com} 长度为 $\ell$ 的细杆沿 $x$ 轴放置，$0\le x\le\ell$，线密度 $\lambda(x)=\lambda_0(1+x/\ell)$。求总质量与质心坐标。
\item\label{exr:mech-restitution} 质量 $m$ 的物体以速率 $u$ 追撞静止的质量 $2m$ 的物体，恢复系数 $e=1/2$。求碰后速度及动能损失占初动能的比例。
\item\label{exr:mech-rocket-gravity} 火箭竖直向上飞行，$m(t)=m_0-qt$，$u_e$ 恒定，重力加速度近似为常数 $g$。求速度随时间的关系，说明公式的时间范围和初始净加速度为正的条件。
\end{enumerate}

\chapter{角动量与转动效应}
\label{chap:mech-angular}

\section{力矩与取矩点}
相同的力作用在门轴附近或远离门轴的位置，转动效果不同。关于指定点 $O$ 的\keyterm{力矩（torque）}定义为 $\boldsymbol\tau_O=\mathbf r\times\mathbf F$，其中 $\mathbf r$ 从 $O$ 指向力的作用点。力矩大小为 $Fd_\perp$，$d_\perp$ 是取矩点到力的作用线的距离；它与力作用点到取矩点的距离一般不同。

质点关于 $O$ 的\keyterm{角动量（angular momentum）}定义为 $\mathbf L_O=\mathbf r\times\mathbf p$。力矩与角动量都依赖取矩点。力矩的量纲与能量相同，但力矩为轴矢量，能量为标量，二者的物理意义不同。

\begin{theorem}[关于惯性系固定点的角动量定理]
若 $O$ 固定于惯性系中，恒质量质点满足 $\dot{\mathbf L}_O=\boldsymbol\tau_O$。对固定成员的质点系，若内力矩总和为零，则
\begin{equation}
 \frac{\mathrm d\mathbf L_O}{\mathrm dt}=\boldsymbol\tau_{O,\mathrm{ext}},\qquad
 \mathbf L_O=\sum_i\mathbf r_i\times m_i\mathbf v_i.
 \label{eq:mech-angular-balance}
\end{equation}
\end{theorem}
\begin{proof}
对单质点，$\mathrm d(\mathbf r\times m\mathbf v)/\mathrm dt=\mathbf v\times m\mathbf v+\mathbf r\times\mathbf F=\boldsymbol\tau_O$。对系统求和，每对内力矩为 $(\mathbf r_i-\mathbf r_j)\times\mathbf F_{ij}$；若内力沿连线，则该项为零。更一般地，只需全部内力矩总和为零。
\end{proof}
总外力为零与总外力矩为零不是同一条件。等大反向但作用线不同的两力构成\keyterm{力偶（couple）}，合力为零而总力矩可不为零。

\section{移动取矩点与质心分解}
设移动点 $A$ 的位置为 $\mathbf r_A(t)$，仍用质点在惯性系中的动量定义
\[
 \mathbf L_A=\sum_i(\mathbf r_i-\mathbf r_A)\times m_i\mathbf v_i.
\]
直接求导并利用总内力矩为零，可得
\begin{equation}
 \dot{\mathbf L}_A=\boldsymbol\tau_{A,\mathrm{ext}}-\mathbf v_A\times\mathbf P.
 \label{eq:mech-moving-origin}
\end{equation}
因此不能随意对移动点沿用固定点形式。若 $A$ 为质心，则 $\mathbf P=M\mathbf V$ 与 $\mathbf v_A=\mathbf V$ 平行，附加项为零，关于质心的角动量定理仍成立，即使质心正在加速。

令 $\mathbf r_i=\mathbf R+\boldsymbol\rho_i$、$\mathbf v_i=\mathbf V+\mathbf v_i'$，利用 $\sum_i m_i\boldsymbol\rho_i=0$ 与 $\sum_i m_i\mathbf v_i'=0$，得到
\begin{equation}
 \mathbf L_O=\mathbf R\times M\mathbf V+\mathbf L_C,\qquad
 \mathbf L_C=\sum_i\boldsymbol\rho_i\times m_i\mathbf v_i'.
 \label{eq:mech-angular-decomposition}
\end{equation}
第一项描述质心绕 $O$ 的运动，第二项描述系统相对质心的内部运动。关于其他移动点若改用相对速度定义角动量，修正公式又不同，使用前必须说明定义。

\section{有心力与面积定律}
\keyterm{有心力（central force）}始终沿质点与固定力心之间的连线，形式为 $\mathbf F=F_r(r,t)\mathbf e_r$。其关于力心的力矩为零，所以 $\mathbf L$ 为常矢量。若 $\mathbf L\ne0$，由 $\mathbf r\cdot\mathbf L=0$，轨迹位于垂直于 $\mathbf L$ 的固定平面内。在此平面内取极坐标，
\[
 L=mr^2\dot\theta,\qquad \frac{\mathrm dA}{\mathrm dt}=\frac12r^2\dot\theta=\frac{L}{2m}.
\]
这里选择角度增加方向使 $L>0$，$A$ 是扫过面积。面积速度恒定来自零力矩，与力是否平方反比无关；轨道为圆锥曲线则需引力的具体形式。若 $\mathbf L=0$，运动为径向的退化情形。

\begin{example}[穿孔细绳收缩圆周运动]
质点在光滑水平桌面上运动，由穿过固定小孔的绳牵引。开始时半径为 $r_1$、切向速率为 $v_1$；收绳后在半径 $r_2$ 处，径向速度再次为零。求末速和拉力总功。
\end{example}
\begin{solve}
拉力关于孔的力矩为零，所以 $mr_1v_1=mr_2v_2$，有 $v_2=r_1v_1/r_2$。由动能定理，拉力功为 $W=\tfrac12mv_1^2(r_1^2/r_2^2-1)$。当收绳时 $r_2<r_1$，拉力做正功，机械能增加；零力矩只保证角动量守恒，不能推出拉力不做功。
\end{solve}

\section{角冲量与守恒的选取}
积分式~\eqref{eq:mech-angular-balance}，得 $\Delta\mathbf L_O=\int\boldsymbol\tau_{O,\mathrm{ext}}\,\mathrm dt$，右端称\keyterm{角冲量（angular impulse）}。碰撞时若支点施加未知冲量，可选支点为取矩点，使支点冲量的力矩为零；但必须检查其他外角冲量是否可忽略。能否消去未知力只是选点的便利，不能代替固定点或修正项条件。

\section{习题}
\begin{enumerate}[label=\thechapter.\arabic*,ref=\thechapter.\arabic*]
\item\label{exr:mech-straight-L} 质点沿直线 $y=b$ 以恒速度 $v\mathbf e_x$ 运动。求关于原点的角动量，并说明没有绕原点作圆周运动是否意味着角动量为零。
\item\label{exr:mech-couple} 两力分别为作用于 $(a,0,0)$ 的 $F\mathbf e_y$ 和作用于 $(-a,0,0)$ 的 $-F\mathbf e_y$。求合力、关于原点的总力矩，并证明其总力矩不随取矩点平移而变。
\item\label{exr:mech-stick-impact} 长 $\ell$、质量 $M$ 的均匀细杆绕一端的固定水平轴在竖直面内转动，初始竖直向下静止。质量 $m$、水平速度 $u$ 的弹丸击中杆的下端并嵌入。已知杆关于轴的转动惯量为 $M\ell^2/3$，求碰后瞬时角速度并解释适用的守恒律。
\end{enumerate}

\chapter{刚体运动与静力平衡}
\label{chap:mech-rigid}

\section{刚体的平动与转动}
刚体任意两点间距离不变，使各点速度并非独立。相对于质心 $C$，刚体上质点的速度可写为
\begin{equation}
 \mathbf v_i=\mathbf V_C+\boldsymbol\omega\times\boldsymbol\rho_i,
 \label{eq:mech-rigid-velocity}
\end{equation}
其中 $\boldsymbol\omega$ 为\keyterm{角速度（angular velocity）}，$\boldsymbol\rho_i$ 从质心指向该点。固定轴转动时，$\boldsymbol\omega$ 沿转轴，垂直距离为 $r_\perp$ 的点速率为 $|\omega|r_\perp$。角速度是一瞬间描述转动的矢量，而有限转动的先后顺序一般不可交换，不能把任意有限转角当作普通矢量相加。

\section{转动惯量与平行轴定理}
\begin{definition}[关于指定轴的转动惯量]
质量分布关于指定轴的\keyterm{转动惯量（moment of inertia）}为
\begin{equation}
 I=\sum_i m_i r_{i\perp}^2=\int r_\perp^2\,\mathrm dm.
 \label{eq:mech-inertia}
\end{equation}
它依赖轴的位置和方向，单位为 $\mathrm{kg\,m^2}$。
\end{definition}
对均匀细杆，垂直于杆且通过中点的轴有 $I_C=(M/\ell)\int_{-\ell/2}^{\ell/2}x^2\,\mathrm dx=M\ell^2/12$。对均匀薄圆盘，绕中心垂直轴将其分成圆环积分，$I_C=\int_0^R r^2(2Mr/R^2)\,\mathrm dr=MR^2/2$。薄圆环绕对称轴则有 $I=MR^2$，说明质量离轴越远，转动惯量越大。

\begin{theorem}[平行轴定理]
设两条轴平行，其中一条通过质心，另一条与它相距 $d$，则
\begin{equation}
 I=I_C+Md^2.
 \label{eq:mech-parallel-axis}
\end{equation}
\end{theorem}
\begin{proof}
在垂直于轴的平面内写 $\mathbf r_\perp=\boldsymbol\rho_\perp+\mathbf d$，展开平方后积分。交叉项为 $2\mathbf d\cdot\int\boldsymbol\rho_\perp\,\mathrm dm=0$，其余两项即为结论。只有其中一条轴通过质心时，交叉项才由这个理由消失。
\end{proof}

\section{固定轴动力学与平面运动}
对刚体绕惯性系中固定的轴转动，轴向角动量为 $L_z=I_z\omega$；即使完整的 $\mathbf L$ 不平行于 $\boldsymbol\omega$，轴向关系仍成立。因此
\begin{equation}
 \tau_z=I_z\dot\omega,\qquad K_{\mathrm{rot}}=\tfrac12I_z\omega^2,
 \qquad P=\tau_z\omega.
 \label{eq:mech-fixed-axis}
\end{equation}
这里的 $I_z$ 恒定，$\tau_z$ 是全部外力矩的轴向分量。平面刚体运动则分解为质心平动与绕质心转动，令角加速度 $\alpha=\dot\omega$，得到
\[
 M\mathbf A_C=\mathbf F_{\mathrm{ext}},\qquad I_C\alpha=\tau_{C,z},\qquad
 K=\tfrac12MV_C^2+\tfrac12I_C\omega^2.
\]
这些式子是相互补充的动力学方程，必须结合几何约束使用。

\begin{example}[具有转动惯量的滑轮]
两物体质量为 $m_2>m_1$，滑轮半径 $R$、绕固定轴的转动惯量 $I$，轴无摩擦，轻绳不可伸长且与滑轮不打滑。求加速度。
\end{example}
\begin{solve}
分别有 $m_2g-T_2=m_2a$、$T_1-m_1g=m_1a$。滑轮受力矩 $(T_2-T_1)R=I\alpha$，不打滑给出 $a=R\alpha$。三式相加，得 $a=(m_2-m_1)g/(m_1+m_2+I/R^2)$。$I/R^2$ 具有质量量纲，反映驱动滑轮转动所需的额外惯性；当 $I\to0$ 时才有 $T_1=T_2$。
\end{solve}

\section{纯滚动与静摩擦}
\keyterm{纯滚动（rolling without slipping）}要求接触点相对接触面瞬时速度为零。在倾角为 $\theta_s$ 的固定直斜面上滚动的轴对称物体满足 $v_C=R\omega$、$a_C=R\dot\omega$，其中正方向取为向下滚动。令 $I_C=\beta MR^2$，假设静摩擦指向斜面上方，方程为
\[
 Mg\sin\theta_s-f=Ma_C,\qquad fR=I_C\frac{a_C}{R}.
\]
解得
\begin{equation}
 a_C=\frac{g\sin\theta_s}{1+\beta},\qquad
 f=\frac{\beta}{1+\beta}Mg\sin\theta_s.
 \label{eq:mech-rolling}
\end{equation}
成立还需 $f\le\mu_sMg\cos\theta_s$，即 $\mu_s\ge\beta\tan\theta_s/(1+\beta)$。若不满足，须取消纯滚动约束并用滑动摩擦模型重新求解。

理想刚体在固定表面纯滚动时，接触点瞬时速度为零，静摩擦的实际功率为零；然而它分别改变质心平动动能和转动动能。把 $f$ 乘质心位移当作摩擦的全部实际功会漏掉转动贡献。滚动阻力、材料形变和移动接触面会改变这一结论。

\section{静力平衡与支承可行性}
刚体静止平衡要求合外力为零，且关于任意一点的合外力矩为零。若已满足合力为零，力矩平衡不随取矩点选择而变。平衡方程给出力的可能值，还需检验静摩擦界限、支持力非负和结构接触等条件；代数解不必对应物理可实现的状态。

\begin{example}[靠在光滑墙上的梯子]
长 $\ell$、质量 $M$ 的均匀梯子与水平面成角 $\alpha$，墙面光滑，地面静摩擦系数为 $\mu_s$。求静止的条件。
\end{example}
\begin{solve}
地面支持力为 $N=Mg$，水平摩擦力大小等于墙的支持力 $N_w$。关于梯脚取矩，$N_w\ell\sin\alpha=Mg(\ell/2)\cos\alpha$，所以 $N_w=(Mg/2)\cot\alpha$。必须满足 $\tfrac12\cot\alpha\le\mu_s$。角度越小，所需摩擦越大；这揭示接近平放的梯子为何难以仅靠地面摩擦维持平衡。
\end{solve}

\section{三维转动与陀螺：选学}
一般三维转动中，绕质心的角动量由\keyterm{惯性张量（inertia tensor）}给出：
\[
 \mathbf L_C=\mathsf I_C\boldsymbol\omega,\qquad
 (\mathsf I_C)_{ab}=\int(\rho^2\delta_{ab}-\rho_a\rho_b)\,\mathrm dm,
\]
其中 $a,b$ 为三个坐标指标，$\delta_{ab}$ 是 Kronecker 符号。实对称矩阵可选正交本征基，对应\keyterm{主轴（principal axes）}。只有绕主轴转动等特殊情形，才有 $\mathbf L_C\parallel\boldsymbol\omega$。在随体主轴系中，角动量方程包含基矢转动项 $\dot{\mathbf L}_{\mathrm{body}}+\boldsymbol\omega\times\mathbf L=\boldsymbol\tau$。

快速自转、缓慢稳定进动的轴对称陀螺可作近似分析。若固定支点到质心距离为 $d$，对称轴与竖直成角 $\theta$，且角动量近似为沿对称轴的 $I_s\omega_s$，则重力矩大小为 $Mgd\sin\theta$。\keyterm{进动（precession）}以角速度 $\Omega_p$ 改变角动量方向，故 $\Omega_p L\sin\theta\approx Mgd\sin\theta$，得到 $\Omega_p\approx Mgd/(I_s\omega_s)$。这是忽略章动和进动自身角动量的近似，不是任意初态、任意自转速度下的精确公式。

\section{习题}
\begin{enumerate}[label=\thechapter.\arabic*,ref=\thechapter.\arabic*]
\item\label{exr:mech-rod-I} 从积分定义求均匀细杆绕端点垂直轴的转动惯量，再用平行轴定理核查。
\item\label{exr:mech-disc-roll} 均匀实心圆柱从静止沿固定斜面纯滚动，下降竖直高度 $h$。求质心速率及维持纯滚动的最小静摩擦系数。
\item\label{exr:mech-physical-pendulum} 刚体绕固定水平轴摆动，质量 $M$，质心到轴距离 $d>0$，绕轴转动惯量 $I_O$。求稳定平衡附近的小振动周期，并求均匀细杆绕端点摆动的周期。
\end{enumerate}

\chapter{万有引力与有心力轨道}
\label{chap:mech-gravity}

\section{引力场、势能与球对称分布}
牛顿万有引力模型中，两质点相距 $r>0$，相互吸引力大小为 $Gm_1m_2/r^2$，$G$ 为引力常量。取两体相对位置 $\mathbf r=\mathbf r_1-\mathbf r_2$，则作用于物体 $1$ 的引力为 $-Gm_1m_2\mathbf r/r^3$。以无限远势能为零，积分得
\begin{equation}
 U(r)=-\frac{Gm_1m_2}{r}.
 \label{eq:mech-gravity-U}
\end{equation}
负号表达将两物体从有限间距缓慢分离到无限远需要输入正功，不表示能量的绝对零点具有特殊可测意义。

球对称物体外部的引力可等效集中于球心，这一结论依赖平方反比规律。以半径 $R$、均匀面密度 $\sigma$ 的薄球壳说明：对离球心 $r>0$ 的试探点，单位质量势为
\[
 \Phi(r)=-2\pi G\sigma R^2\int_{-1}^1\frac{\mathrm du}{\sqrt{r^2+R^2-2rRu}}
 =-\frac{GM}{2rR}\bigl(r+R-|r-R|\bigr),
\]
其中 $M=4\pi R^2\sigma$，变量 $u$ 为极角余弦。因此壳外 $\Phi=-GM/r$，壳内 $\Phi=-GM/R$ 为常数，壳内引力为零。将任意球对称密度分布分成薄壳叠加，就得到壳外等效结论；均匀实心球内部则有 $g(r)=GMr/R^3$。球心处取连续极限，不直接使用含 $1/r$ 的中间表达式。

\section{圆轨道与逃逸条件}
当中心天体质量 $M$ 远大于卫星质量 $m$，忽略中心天体反冲，可把力心视为固定。半径为 $r$ 的圆轨道满足
\[
 v_c=\sqrt{\frac{GM}{r}},\qquad
 T=2\pi\sqrt{\frac{r^3}{GM}},\qquad E=-\frac{GMm}{2r}.
\]
从半径 $r$ 处摆脱引力束缚并在无限远速率趋于零，最低能量要求 $E=0$，得到\keyterm{逃逸速率（escape speed）} $v_e=\sqrt{2GM/r}$。这忽略大气阻力、自转、其他天体及推进过程；对于有实体表面的天体，发射方向还需避免轨迹与天体相交。能量非负描述不受束缚的轨道性质，不是任意初始方向都能实际离开表面。

\section{两体问题与有效势}
对于孤立两体，质心匀速运动。相对坐标 $\mathbf r=\mathbf r_1-\mathbf r_2$ 满足
\[
 \mu\ddot{\mathbf r}=-\frac{k}{r^3}\mathbf r,\qquad
 \mu=\frac{m_1m_2}{m_1+m_2},\qquad k=Gm_1m_2.
\]
因此两体内部运动化为质量 $\mu$ 的等效质点运动。相对运动能量记为 $E_{\mathrm{rel}}$，角动量大小为 $L=\mu r^2\dot\theta$。对 $L>0$，有
\begin{equation}
 E_{\mathrm{rel}}=\tfrac12\mu\dot r^2+U_{\mathrm{eff}}(r),\qquad
 U_{\mathrm{eff}}(r)=\frac{L^2}{2\mu r^2}-\frac{k}{r}.
 \label{eq:mech-effective-potential}
\end{equation}
\keyterm{有效势（effective potential）}中的第一项来自固定角动量下的切向动能，不是另加了一种真实相互作用。极小值满足 $r_c=L^2/(\mu k)$，对应圆轨道；非圆束缚轨道的径向距离在两个转折点之间变化。

\section{轨道方程与开普勒定律}
令 $u=1/r$，撇号表示对 $\theta$ 求导。由角动量守恒，$\dot r=-(L/\mu)u'$，$\ddot r=-(L^2/\mu^2)u^2u''$，且 $r\dot\theta^2=(L^2/\mu^2)u^3$。代入径向牛顿方程，得到
\[
 u''+u=\frac{\mu k}{L^2}.
\]
选择近心点为 $\theta=0$，解为
\begin{equation}
 r(\theta)=\frac{p}{1+e\cos\theta},\qquad
 p=\frac{L^2}{\mu k},\qquad
 e=\sqrt{1+\frac{2E_{\mathrm{rel}}L^2}{\mu k^2}}.
 \label{eq:mech-orbit}
\end{equation}
最后一个关系可由近心点 $\dot r=0$ 代入式~\eqref{eq:mech-effective-potential} 得到。$e$ 为\keyterm{离心率（eccentricity）}：$0\le e<1$ 为椭圆，$e=1$ 为抛物线，$e>1$ 为双曲线；$L=0$ 的径向情形须另行处理。

椭圆半长轴 $a=p/(1-e^2)$、半短轴 $b=a\sqrt{1-e^2}$，相对运动能量为 $E_{\mathrm{rel}}=-k/(2a)$。面积速度为 $L/(2\mu)$，一周扫过面积 $\pi ab$，所以
\begin{equation}
 T^2=\frac{4\pi^2a^3}{G(m_1+m_2)}.
 \label{eq:mech-kepler-third}
\end{equation}
这给出两体问题的开普勒第三定律精确牛顿形式；$a$ 是相对轨道半长轴。当一个质量远大于另一个时才可近似把质量和替换为中心天体质量。第一定律的椭圆轨道与第二定律的面积定律分别来自平方反比力和角动量守恒，逻辑依据不同。

\begin{example}[近地点和远地点的速度比]
忽略卫星对中心天体的反冲，已知椭圆轨道最近、最远距离为 $r_p,r_a$。求两处速率及其比例。
\end{example}
\begin{solve}
两处径向速度为零，角动量守恒给出 $r_pv_p=r_av_a$。又 $a=(r_p+r_a)/2$，由能量 $v^2/2-GM/r=-GM/(2a)$，得
\[
 v_p^2=\frac{2GMr_a}{r_p(r_p+r_a)},\qquad
 v_a^2=\frac{2GMr_p}{r_a(r_p+r_a)},\qquad \frac{v_p}{v_a}=\frac{r_a}{r_p}.
\]
近地点速度较大并不是因为那里有切向引力分量；在近地点瞬间引力为径向力，速度较大来自此前运动中势能向动能的转化。
\end{solve}

\section{习题}
\begin{enumerate}[label=\thechapter.\arabic*,ref=\thechapter.\arabic*]
\item\label{exr:mech-satellite} 同一中心天体的两条圆轨道半径之比为 $4$。求轨道速率和周期之比，并说明比较时采用的质量近似。
\item\label{exr:mech-uniform-earth} 假设天体是半径 $R$、质量 $M$ 的均匀球体，忽略自转，在穿过球心的细直隧道内释放试验质点，忽略隧道挖去的质量及阻力。求运动方程与周期。
\item\label{exr:mech-orbit-energy} 两质量相等、各为 $m$ 的天体作圆形相对运动，间距为 $d$。求各自相对质心的速率、周期及系统在质心系中的总能量。
\end{enumerate}

\chapter{非惯性参考系}
\label{chap:mech-frames}

\section{平动加速系与惯性力}
在加速列车中，悬挂物的平衡方向会偏离竖直。若参考系 $S'$ 的轴与惯性系平行，原点加速度为 $\mathbf A$，由 $\mathbf a=\mathbf A+\mathbf a'$ 得
\[
 m\mathbf a'=\mathbf F_{\mathrm{real}}-m\mathbf A.
\]
将 $-m\mathbf A$ 称为\keyterm{惯性力（inertial force）}。它是改写运动方程产生的参考系项，不对应另一个物体的相互作用，也不要求为其单独寻找作用反作用力。

\begin{example}[加速车厢中的摆]
车厢以恒定加速度 $A>0$ 向右运动，长度 $\ell$ 的单摆相对车厢保持静止，求摆线倾角及稳定平衡附近的小振动周期。
\end{example}
\begin{solve}
车厢系中的有效重力为 $\mathbf g_{\mathrm{eff}}=-A\mathbf e_x-g\mathbf e_y$，摆球向左偏，摆线与竖直向下的夹角满足 $\tan\theta=A/g$。有效重力大小为 $\sqrt{g^2+A^2}$，因此小振动周期为 $2\pi\sqrt{\ell/\sqrt{g^2+A^2}}$。加速度恒定且车厢不转动是直接使用固定有效重力的条件。
\end{solve}

\section{转动坐标基的求导法则}
设 $S'$ 相对惯性系以角速度 $\boldsymbol\Omega(t)$ 转动。对任意矢量 $\mathbf B$，其分量变化与基矢变化相加，给出
\begin{equation}
 \left(\frac{\mathrm d\mathbf B}{\mathrm dt}\right)_S
 =\left(\frac{\mathrm d\mathbf B}{\mathrm dt}\right)_{S'}+\boldsymbol\Omega\times\mathbf B.
 \label{eq:mech-transport}
\end{equation}
这里两侧描述同一个几何矢量；下标区分在何种坐标基中求导。若转动系原点在惯性系中的位置为 $\mathbf R(t)$，质点相对位置为 $\mathbf r'$，一次求导得 $\mathbf v=\dot{\mathbf R}+\mathbf v'+\boldsymbol\Omega\times\mathbf r'$。再次求导并对 $\mathbf v'$ 与 $\boldsymbol\Omega\times\mathbf r'$ 分别使用式~\eqref{eq:mech-transport}，得到
\begin{equation}
 \mathbf a=\mathbf A+\mathbf a'+2\boldsymbol\Omega\times\mathbf v'
 +\dot{\boldsymbol\Omega}\times\mathbf r'
 +\boldsymbol\Omega\times(\boldsymbol\Omega\times\mathbf r').
 \label{eq:mech-rotating-acceleration}
\end{equation}
系数 $2$ 来自两个不同求导项，不能从单纯的圆周向心加速度猜出。

\section{离心力、科里奥利力与欧拉力}
将式~\eqref{eq:mech-rotating-acceleration} 代入牛顿方程，移项得
\begin{align}
 m\mathbf a'={}&\mathbf F_{\mathrm{real}}-m\mathbf A
 -2m\boldsymbol\Omega\times\mathbf v'
 -m\dot{\boldsymbol\Omega}\times\mathbf r'\notag\\
 &-m\boldsymbol\Omega\times(\boldsymbol\Omega\times\mathbf r').
 \label{eq:mech-rotating-newton}
\end{align}
其中 $-2m\boldsymbol\Omega\times\mathbf v'$ 是\keyterm{科里奥利力（Coriolis force）}，$-m\dot{\boldsymbol\Omega}\times\mathbf r'$ 是\keyterm{欧拉力（Euler force）}，最后一项是\keyterm{离心力（centrifugal force）}。离心力指向远离转轴的方向，其大小为 $m\Omega^2r_\perp$。

科里奥利力总与相对速度垂直，故在转动系中瞬时功率为零，但能改变运动方向。恒定角速度且原点固定时，离心力有势能 $U_{\mathrm{cf}}=-m|\boldsymbol\Omega\times\mathbf r'|^2/2$；若真实势能在转动系中也不显含时间，则 $m(v')^2/2+U+U_{\mathrm{cf}}$ 守恒。这一量不同于惯性系中的机械能。

\begin{example}[旋转直槽中的珠]
光滑径向直槽在水平面内绕固定端以恒定角速度 $\Omega>0$ 转动。质量 $m$ 的小珠可沿槽滑动，初始 $r(0)=r_0>0$、$\dot r(0)=0$。忽略空气阻力，求 $r(t)$ 及槽对珠的切向力。
\end{example}
\begin{solve}
在随槽转动系中，真实径向力为零，故 $\ddot r=\Omega^2r$，得到 $r=r_0\cosh(\Omega t)$。珠的相对切向加速度为零，槽必须抵消科里奥利力，所以其切向力为 $N_\theta=2m\Omega\dot r=2m\Omega^2r_0\sinh(\Omega t)$，正方向为槽的转动方向。该力在惯性系中做功，保持转速的驱动机构供能。解仅在珠尚未到达槽端、驱动能够保持规定角速度的阶段有效。
\end{solve}

\section{地球参考系中的量级判断}
地球自转带来的离心项可部分并入当地测得的有效重力，科里奥利项则与运动速度有关。若特征速率为 $U$、水平尺度为 $L$，惯性加速度的量级为 $U^2/L$，科里奥利加速度量级为 $2\Omega U$，二者之比约为 $2\Omega L/U$，具体投影还取决于纬度和方向。大尺度、长时间运动中即使 $\Omega$ 数值较小，也可能不可忽略。用测得的当地 $g$ 后再完整重复添加同一离心修正，会造成重复计入。

\section{习题}
\begin{enumerate}[label=\thechapter.\arabic*,ref=\thechapter.\arabic*]
\item\label{exr:mech-fictitious} 匀速转盘上的物体相对转盘静止。分别在惯性系和转盘系中说明维持圆周运动的合力关系，解释向心力与离心力是否是一对作用反作用力。
\item\label{exr:mech-turntable} 物块放在水平转盘上，距轴 $r$，静摩擦系数为 $\mu_s$，转盘以恒角速度 $\Omega$ 转动。求不滑动条件。
\item\label{exr:mech-coriolis} 在水平面内取 $\boldsymbol\Omega=\Omega\mathbf e_z$，珠的相对速度为 $u\mathbf e_x$。求科里奥利力方向，并证明它对任意相对速度的瞬时功率为零。
\end{enumerate}

\chapter{振动、阻尼与共振}
\label{chap:mech-oscillation}

\section{稳定平衡附近的简谐运动}
许多系统在稳定平衡附近具有相同的线性化方程。设势能在 $x=0$ 处有严格非退化极小值，$U'(0)=0$、$U''(0)=k>0$，则小位移时 $F=-U'(x)\approx-kx$。保留一阶项得到
\begin{equation}
 \ddot x+\omega_0^2x=0,\qquad \omega_0=\sqrt{k/m},\qquad
 x(t)=A\cos\omega_0t+B\sin\omega_0t.
 \label{eq:mech-shm}
\end{equation}
这种运动称为\keyterm{简谐运动（simple harmonic motion）}。初值 $x_0,v_0$ 给出 $A=x_0$、$B=v_0/\omega_0$，振幅为 $\sqrt{x_0^2+(v_0/\omega_0)^2}$，周期 $T_0=2\pi/\omega_0$。相位必须同时满足位置和速度条件，不能只用一个反余弦确定象限。

弹簧振子中 $E=m\dot x^2/2+kx^2/2$ 恒定，动能与势能交替转化。相空间中以 $x$ 和 $p=m\dot x$ 为坐标，固定正能量的轨迹是椭圆。它将一个时刻的位置与速度同时记录，比仅画空间轨迹更能描述动力学状态。

\section{单摆与小角度误差}
单摆的精确约束方程为 $\ddot\theta+(g/\ell)\sin\theta=0$。小角度下才得到 $\omega_0=\sqrt{g/\ell}$。对于最大摆角 $0<\theta_0<\pi$ 的无阻力往复摆动，能量守恒给出
\[
 \dot\theta^2=\frac{2g}{\ell}(\cos\theta-\cos\theta_0).
\]
利用 $\sin(\theta/2)=\sin(\theta_0/2)\sin\phi$，四分之一周期积分化为
\begin{equation}
 T=4\sqrt{\frac\ell g}\int_0^{\pi/2}
 \frac{\mathrm d\phi}{\sqrt{1-\sin^2(\theta_0/2)\sin^2\phi}}.
 \label{eq:mech-pendulum-period}
\end{equation}
小摆角时展开被积函数，得到 $T=T_0[1+\theta_0^2/16+O(\theta_0^4)]$。展开在摆幅远离 $\pi$ 的足够小邻域内可一致控制；当 $\theta_0\to\pi$，周期趋于无穷，小角度公式完全失去精度。周期与振幅无关是线性模型的性质。

\section{阻尼振动}
在线性回复力和线性阻力模型下，令 $2\gamma=b/m$，则
\begin{equation}
 \ddot x+2\gamma\dot x+\omega_0^2x=0,\qquad \gamma\ge0.
 \label{eq:mech-damped}
\end{equation}
将试探解 $e^{st}$ 代入，特征根为 $s=-\gamma\pm\sqrt{\gamma^2-\omega_0^2}$。当 $0<\gamma<\omega_0$ 时为\keyterm{欠阻尼（underdamping）}，解为 $x=e^{-\gamma t}(A\cos\omega_dt+B\sin\omega_dt)$，其中 $\omega_d=\sqrt{\omega_0^2-\gamma^2}$。当 $\gamma=\omega_0$ 时为\keyterm{临界阻尼（critical damping）}，解为 $(A+Bt)e^{-\omega_0t}$；当 $\gamma>\omega_0$ 时为\keyterm{过阻尼（overdamping）}，解为两个衰减指数的线性组合。

三类解都必须结合初值解释，不能仅用“阻尼越大越快停止”概括。过大的阻尼会使慢衰减根接近零，回到平衡反而缓慢。机械能满足 $\dot E=-b\dot x^2\le0$；欠阻尼的振幅包络按 $e^{-\gamma t}$ 衰减，能量只在弱阻尼的周期平均意义下近似按 $e^{-2\gamma t}$ 衰减，瞬时能量还带有周期调制。

\section{受迫振动与频率响应}
若外力为 $F_0\cos\omega t$，取 $F_0>0$、$\omega>0$，则
\[
 \ddot x+2\gamma\dot x+\omega_0^2x=(F_0/m)\cos\omega t.
\]
当 $\gamma>0$ 时齐次瞬态逐渐衰减，稳态解为 $x_{\mathrm{st}}=A(\omega)\cos(\omega t-\delta)$。代入并比较正余弦系数，得
\begin{equation}
 A(\omega)=\frac{F_0/m}{\sqrt{(\omega_0^2-\omega^2)^2+4\gamma^2\omega^2}},
 \qquad \delta=\operatorname{atan2}(2\gamma\omega,\omega_0^2-\omega^2).
 \label{eq:mech-forced-amplitude}
\end{equation}
这里 $\operatorname{atan2}(y,x)$ 按两个参数确定象限，故 $0<\delta<\pi$。\keyterm{共振（resonance）}需要明确考察的响应量：位移振幅在 $\gamma<\omega_0/\sqrt2$ 时于 $\omega=\sqrt{\omega_0^2-2\gamma^2}$ 达峰，平均吸收功率 $\overline P=b\omega^2A^2/2$ 则在恒定驱动力振幅下于 $\omega=\omega_0$ 达峰。

无阻尼且驱动频率恰为 $\omega_0$ 时，不存在有限振幅的上述周期稳态；一个特解为 $x_p=F_0t\sin(\omega_0t)/(2m\omega_0)$。振幅线性增长只在理想线性模型仍有效的阶段可信，实际系统中的阻尼、非线性和驱动限制会改变长期行为。

\begin{example}[由初值确定欠阻尼解]
欠阻尼振子满足 $x(0)=x_0$、$\dot x(0)=0$。求位移。
\end{example}
\begin{solve}
由 $A=x_0$ 及 $-\gamma A+\omega_dB=0$，得
\[
 x(t)=x_0e^{-\gamma t}\left(\cos\omega_dt+\frac\gamma{\omega_d}\sin\omega_dt\right).
\]
只写 $x_0e^{-\gamma t}\cos\omega_dt$ 会把初速度误写为 $-\gamma x_0$。初值检查比套用振幅包络更可靠。
\end{solve}

\section{耦合振动与简正模}
两质量相等为 $m$ 的物块分别由劲度系数 $k$ 的弹簧连至两侧固定墙，二者间以劲度系数 $\kappa$ 的弹簧相连，所有位移从平衡位置量起。方程为
\[
 m\ddot x_1=-(k+\kappa)x_1+\kappa x_2,\qquad
 m\ddot x_2=\kappa x_1-(k+\kappa)x_2.
\]
相加、相减，定义 $q_+=x_1+x_2$、$q_-=x_1-x_2$，得到相互独立的方程 $\ddot q_++(k/m)q_+=0$、$\ddot q_-+[(k+2\kappa)/m]q_-=0$。这两种保持固定振幅比和相位关系的运动称为\keyterm{简正模（normal modes）}。一般初始状态是两模叠加；耦合振动并不意味着每个物块只有一个固有频率。

\section{习题}
\begin{enumerate}[label=\thechapter.\arabic*,ref=\thechapter.\arabic*]
\item\label{exr:mech-spring-vertical} 质量 $m$ 的物体挂在竖直弹簧下，劲度系数 $k$。从自然长度位置静止释放，求平衡伸长、最大伸长及振动周期。
\item\label{exr:mech-log-decrement} 欠阻尼振动的相邻同向位移峰值之比为 $q>1$，两峰间隔为 $T_d$。求 $\gamma$、$\omega_d$ 和 $\omega_0$。
\item\label{exr:mech-modes} 对本章两个耦合振子的模型，初值为 $x_1(0)=a$、$x_2(0)=0$，初速度均为零。求 $x_1(t),x_2(t)$，并说明拍现象对应的频率条件。
\end{enumerate}

\chapter{机械波与连续振动}
\label{chap:mech-waves}

\section{从离散质点到弦的波动方程}
波是扰动通过相互作用传播的过程，介质质点通常只在平衡位置附近振动。设均匀细弦沿 $x$ 轴张紧，张力大小为常数 $\mathcal T$，线密度为 $\lambda$，横向位移为 $y(x,t)$。假设斜率 $|\partial_xy|\ll1$，忽略弦的弯曲刚度和阻尼，且纵向运动对张力的影响为高阶小量。微元 $\Delta x$ 的横向合力为 $\mathcal T[y_x(x+\Delta x,t)-y_x(x,t)]$，所以在足够光滑的连续极限下
\begin{equation}
 \frac{\partial^2y}{\partial t^2}=c_w^2\frac{\partial^2y}{\partial x^2},\qquad
 c_w=\sqrt{\frac{\mathcal T}{\lambda}}.
 \label{eq:mech-wave}
\end{equation}
这是\keyterm{波动方程（wave equation）}。波速 $c_w$ 来自介质的回复能力和惯性，其量纲为速度；它不是弦上质点的横向速度 $\partial_ty$。

\section{行波、初值与叠加}
对无限长、均匀且满足线性模型的弦，足够光滑的函数 $f(x-c_wt)$ 和 $g(x+c_wt)$ 分别描述向右和向左传播的\keyterm{行波（traveling wave）}。其和仍满足方程，这叫\keyterm{叠加原理（superposition principle）}。给定 $y(x,0)=y_0(x)$、$y_t(x,0)=v_0(x)$，可解出
\begin{equation}
 y(x,t)=\frac{y_0(x-c_wt)+y_0(x+c_wt)}2
 +\frac1{2c_w}\int_{x-c_wt}^{x+c_wt}v_0(\xi)\,\mathrm d\xi.
 \label{eq:mech-dalembert}
\end{equation}
这一达朗贝尔公式适用于无限线上的经典解，例如 $y_0$ 二阶连续可微、$v_0$ 一阶连续可微时。有限弦还必须满足端点边界条件；不能把无限弦公式不加修改地使用到反射发生以后。

正弦行波写成 $y=A\cos(k_wx-\omega t+\phi)$，其中 $k_w=2\pi/\Lambda$ 为波数，$\Lambda$ 为波长，$\omega=2\pi f$ 为角频率。代入方程得 $\omega=c_wk_w$，故 $c_w=f\Lambda$。线性、均匀、无色散是所有这类波包以同一速度传播的条件；更复杂介质可能具有依赖波数的传播速度。

\section{反射、驻波与边界条件}
固定端满足 $y=0$，理想自由端满足 $\partial_xy=0$。入射波与反射波在固定端的位移相消，在自由端的斜率相消。对长度 $\ell$、两端固定的弦，可分离变量为 $y=X(x)Q(t)$，得到 $X''+k_w^2X=0$。由 $X(0)=X(\ell)=0$，非零解要求
\begin{equation}
 k_n=\frac{n\pi}{\ell},\qquad \omega_n=\frac{n\pi c_w}{\ell},\qquad
 y_n=A_n\sin\frac{n\pi x}{\ell}\cos(\omega_nt+\phi_n),\quad n=1,2,\ldots.
 \label{eq:mech-string-modes}
\end{equation}
这些模式形成\keyterm{驻波（standing wave）}，波节位置不随时间移动。一般有限能量初态可作正弦级数展开；逐项二阶求导需要更强的光滑性和收敛条件，不能仅凭形式级数就保证经典解。

\begin{example}[拨弦初态的模态展开]
两端固定的弦在中点被拉起高度 $a$ 后从静止释放，初形为 $y_0(x)=2ax/\ell$（$0\le x\le\ell/2$）和 $2a(\ell-x)/\ell$（$\ell/2\le x\le\ell$）。求各简正模系数。
\end{example}
\begin{solve}
初速度为零，故取 $y=\sum_{n\ge1}B_n\sin(n\pi x/\ell)\cos\omega_nt$。利用正交性，
\[
 B_n=\frac2\ell\int_0^\ell y_0(x)\sin\frac{n\pi x}{\ell}\,\mathrm dx
 =\frac{8a}{n^2\pi^2}\sin\frac{n\pi}2.
\]
偶数模消失，奇数模系数交替变号。初形在中点有折角，因此应把该展开理解为有限能量解；沿折角传播的特征线不能无条件要求位移二阶连续可微。
\end{solve}

\section{波的能量与功率}
在小斜率近似下，弦单位平衡长度的动能为 $\lambda y_t^2/2$，由于伸长 $\mathrm ds-\mathrm dx\approx y_x^2\mathrm dx/2$，弹性势能增量为 $\mathcal T y_x^2/2$。定义能量密度 $\mathcal E$ 和沿正 $x$ 方向的功率 $\mathcal P$：
\[
 \mathcal E=\tfrac12\lambda y_t^2+\tfrac12\mathcal T y_x^2,\qquad
 \mathcal P=-\mathcal T y_xy_t.
\]
对光滑解求导并利用式~\eqref{eq:mech-wave}，得 $\partial_t\mathcal E+\partial_x\mathcal P=0$。因此一个弦段能量的变化来自两端能量流，而不要求弦段内部每个位置的能量密度恒定。向右正弦行波的平均功率为 $\overline{\mathcal P}=\mathcal T k_w\omega A^2/2=\lambda c_w\omega^2A^2/2$。

\section{声波与多普勒效应入门}
流体中的小振幅声波是压强和密度的纵向扰动。均匀静止介质中的声速满足 $c_s^2=(\partial p/\partial\rho)_s$，下标 $s$ 表示比熵保持不变的绝热小扰动响应；具体数值还需介质的状态方程。它不同于拉紧弦的横波模型。

对一维经典声波，若声源以速度大小 $v_s<c_s$ 朝向接收者运动，接收者以速度大小 $v_o$ 朝向声源运动，二者速度均相对介质测量，则相邻波峰间距为 $(c_s-v_s)/f_s$，接收者遇波峰的相对速率为 $c_s+v_o$，所以
\[
 f_o=f_s\frac{c_s+v_o}{c_s-v_s}.
\]
远离时应按方向改变符号。该\keyterm{多普勒效应（Doppler effect）}公式依赖介质参考系，不能直接用于真空光波。

\section{习题}
\begin{enumerate}[label=\thechapter.\arabic*,ref=\thechapter.\arabic*]
\item\label{exr:mech-wave-data} 弦波 $y=A\cos(k_wx-\omega t)$ 中 $A=2\,\mathrm{mm}$、$k_w=4\pi\,\mathrm{m^{-1}}$、$\omega=20\pi\,\mathrm{s^{-1}}$。求波长、频率、波速及质点横向速度最大值。
\item\label{exr:mech-fixed-free} 长 $\ell$ 的弦一端固定、一端理想自由，波速为 $c_w$。求允许的简正模频率。
\item\label{exr:mech-wave-energy} 对两端固定弦的单模 $y=A\sin(n\pi x/\ell)\cos\omega_nt$，直接积分求总能量，验证它与时间无关。
\end{enumerate}

\chapter{弹性与流体力学基础}
\label{chap:mech-continuum}

\section{连续介质与线弹性}
\keyterm{连续介质（continuum）}模型把密度、位移和速度视为空间中的连续场，适用尺度必须足够大，能够平均掉微观离散结构，又足够小，能分辨宏观变化。质点、刚体和连续介质不是互相排斥的世界观，而是保留不同自由度的建模层次。

均匀细杆沿轴向受拉，原长 $\ell$、截面积 $A$、伸长 $\Delta\ell$。\keyterm{正应力（normal stress）}为 $\sigma=F/A$，\keyterm{应变（strain）}为 $\varepsilon=\Delta\ell/\ell$。小形变线弹性范围内，$\sigma=Y\varepsilon$，$Y$ 为\keyterm{杨氏模量（Young's modulus）}。因此杆的等效劲度系数为 $YA/\ell$，弹性能为
\[
 U=\frac12\frac{YA}{\ell}(\Delta\ell)^2=\frac12Y\varepsilon^2A\ell.
\]
这是材料本构关系，超过线性范围或出现塑性形变后不能继续使用。剪切模量与体积模量分别描述剪切和体积压缩；体积模量定义为 $B=-V\,\mathrm dp/\mathrm dV$，正号表示稳定材料通常越压缩压强越大。

\section{静水压强与浮力}
\keyterm{压强（pressure）}是各向同性压缩应力的大小，静止流体的压强力在面元上沿法向作用。取竖直向上 $z$ 轴，厚度 $\mathrm dz$、水平面积 $A$ 的静止流体微元受力平衡给出
\begin{equation}
 \frac{\mathrm dp}{\mathrm dz}=-\rho g.
 \label{eq:mech-hydrostatic}
\end{equation}
密度恒定时，深度为 $h$ 处 $p=p_0+\rho gh$；密度随深度变化时必须积分，不能将某一点的密度乘总深度。压强是标量，压强力有方向，绝对压强与相对大气的表压也须区分。

浸入流体的物体受到周围压强力。设排开流体占据体积 $V_d$，由散度定理，合力为
\[
 \mathbf F_b=-\oint_{\partial V_d}p\mathbf n\,\mathrm dA
 =-\int_{V_d}\nabla p\,\mathrm dV
 =\int_{V_d}\rho g\mathbf e_z\,\mathrm dV.
\]
这就是\keyterm{阿基米德原理（Archimedes' principle）}：浮力等于被排开流体的重量，方向向上。推导要求包围物体的压力分布可由相应静水场描述；物体与容器底密封贴合时，底面压力条件改变，不能照搬完整浸没物体的公式。

\begin{example}[漂浮物的浸没比例]
密度均匀为 $\rho_b$ 的物体漂浮于密度 $\rho_f>\rho_b$ 的静止液体表面，忽略空气浮力和表面张力，求浸没体积比例。
\end{example}
\begin{solve}
平衡时 $\rho_f gV_d=\rho_bgV$，故 $V_d/V=\rho_b/\rho_f$。这一结果只确定体积比例，不确定物体的朝向；转动稳定性还依赖形状、质心和浮心的位置。
\end{solve}

\section{质量守恒与连续性方程}
设密度为 $\rho(\mathbf x,t)$、流速场为 $\mathbf u(\mathbf x,t)$，取固定控制体 $V$，外法向为 $\mathbf n$。质量守恒为
\[
 \frac{\mathrm d}{\mathrm dt}\int_V\rho\,\mathrm dV=-\oint_{\partial V}\rho\mathbf u\cdot\mathbf n\,\mathrm dA.
\]
若场足够光滑，可将导数移入积分并用散度定理，得到
\begin{equation}
 \partial_t\rho+\nabla\cdot(\rho\mathbf u)=0.
 \label{eq:mech-continuity}
\end{equation}
\keyterm{不可压缩流动（incompressible flow）}满足流体微团沿运动的密度不变。用 $D\rho/Dt=\partial_t\rho+\mathbf u\cdot\nabla\rho$ 表示沿微团轨迹的密度变化率，则不可压缩条件为 $D\rho/Dt=0$；将式~\eqref{eq:mech-continuity} 展开可得 $\nabla\cdot\mathbf u=0$。对于定常、单入口单出口且截面内速度近似均匀的管流，$\rho Av$ 为常数；密度相同才进一步得到 $Av$ 为常数。

\section{欧拉方程与伯努利方程}
流体微团沿轨迹的加速度不是单纯的 $\partial_t\mathbf u$，还包含进入不同空间位置引起的变化。\keyterm{物质导数（material derivative）}为
\[
 \frac{D\mathbf u}{Dt}=\partial_t\mathbf u+(\mathbf u\cdot\nabla)\mathbf u.
\]
在忽略黏性、只有压强力和重力的模型下，微元动量平衡为 $\rho D\mathbf u/Dt=-\nabla p+\rho\mathbf g$，称为理想流体的欧拉方程。对定常、恒密度、无黏性流动，沿流线取切向分量并积分，得到
\begin{equation}
 \frac12v^2+\frac p\rho+gz=C_{\mathrm{streamline}}.
 \label{eq:mech-bernoulli}
\end{equation}
这就是\keyterm{伯努利方程（Bernoulli equation）}的常用形式。常数一般只能在同一条流线上比较；若流动在连通区域内还无旋，常数才可跨流线统一。泵输入、黏性损失、明显可压缩性或非定常性都需要修改模型。

\begin{example}[大水箱的小孔出流]
水箱自由液面到小孔的高度差为 $h$，自由液面面积远大于小孔面积，两处都接触同一大气压。忽略黏性损失，采用准定常近似，求小孔流速。
\end{example}
\begin{solve}
质量守恒表明自由液面下降速度远小于出口速度。沿连接液面和出口的流线应用式~\eqref{eq:mech-bernoulli}，两端压强相消，得到 $v\approx\sqrt{2gh}$。若水位显著下降，$h$ 随时间变化，公式应作为随时更新的准定常关系，而不是整个排空过程中速度恒定的结论。
\end{solve}

\section{黏性模型的适用边界}
\keyterm{黏性（viscosity）}描述流体内部剪切运动造成的动量传输。牛顿流体的简单剪切满足 $\tau=\eta\,\mathrm du/\mathrm dy$，其中 $\eta$ 为动力黏度。尺度为 $L$、速度为 $U$ 的流动可用\keyterm{雷诺数（Reynolds number）} $\mathrm{Re}=\rho UL/\eta$ 比较惯性与黏性项的量级。低雷诺数常允许忽略惯性项；高雷诺数并不意味着所有区域黏性都无关，固壁附近的边界层仍可能控制阻力。层流与湍流的转变取决于几何和扰动，不能给全部流动指定一个统一临界雷诺数。

\section{习题}
\begin{enumerate}[label=\thechapter.\arabic*,ref=\thechapter.\arabic*]
\item\label{exr:mech-rods} 两根长度分别为 $\ell_1,\ell_2$、截面积同为 $A$、杨氏模量为 $Y_1,Y_2$ 的细杆串联受拉。忽略自重，求等效劲度系数。
\item\label{exr:mech-dam} 竖直矩形闸门宽 $b$、高 $h$，顶部与液面齐平，两侧上方均为大气，另一侧无液体。求液体对闸门的净水平力及作用点到液面的深度。
\item\label{exr:mech-pipe} 水平管道两截面面积为 $A_1,A_2$，流体恒密度为 $\rho$，定常无黏性，截面内速度均匀。已知体积流量 $Q$，求 $p_1-p_2$，并判断缩窄处压强如何变化。
\end{enumerate}

\chapter{狭义相对论力学入门}
\label{chap:mech-relativity}

\section{时空测量与基本原理}
牛顿力学假设不同惯性系共享同一时间，伽利略速度叠加据此成立。狭义相对论以不同的时空结构描述高速运动，其基本出发点是：物理规律在所有惯性系中形式相同；真空光速 $c$ 在各惯性系中相同，与光源运动无关。这些是物理原理，并非从牛顿定律推出的数学命题。本章只讨论平直时空，忽略引力对时空的影响。

\keyterm{事件（event）}是发生在某个空间位置和某个时刻的物理过程点，用 $(t,x,y,z)$ 标记。比较两个相隔位置的时钟，需要先约定同步方法；在惯性系中可用往返光信号定义爱因斯坦同步。钟显示的读数、坐标时间和“肉眼何时看见”不能混为一谈，后者还含光传播延迟。

\section{洛伦兹变换}
设 $S'$ 相对 $S$ 沿正 $x$ 方向以恒速 $V$ 运动，$|V|<c$，两系轴平行且原点在零时刻重合。空间、时间均匀性要求变换为线性；$S'$ 原点轨迹 $x=Vt$ 使 $x'=a(x-Vt)$。要求沿 $\pm x$ 方向的光线在两系都满足 $x=\pm ct$、$x'=\pm ct'$，得 $t'=a(t-Vx/c^2)$。互易性与空间各向同性使逆变换的系数相同且 $V$ 反号，复合后要求 $a^2(1-V^2/c^2)=1$。取连续于 $V=0$ 的正根，得到\keyterm{洛伦兹变换（Lorentz transformation）}
\begin{equation}
 x'=\Gamma(x-Vt),\qquad t'=\Gamma\left(t-\frac{Vx}{c^2}\right),\qquad
 y'=y,\quad z'=z,\quad \Gamma=\frac1{\sqrt{1-V^2/c^2}}.
 \label{eq:mech-lorentz}
\end{equation}
直接展开可检验 $c^2\Delta t'^2-\Delta x'^2-\Delta y'^2-\Delta z'^2=c^2\Delta t^2-\Delta x^2-\Delta y^2-\Delta z^2$。该不变量称为\keyterm{时空间隔（spacetime interval）}，本书采用时间项为正的符号约定。

\section{同时性、固有时与长度收缩}
在 $S$ 中同时的两个事件有 $\Delta t=0$，若 $\Delta x\ne0$，通常 $\Delta t'=-\Gamma V\Delta x/c^2\ne0$。因此同时性与参考系有关。\keyterm{固有时（proper time）}是随物体运动的理想钟沿自身轨迹记录的时间，对于速率 $v<c$ 的钟有
\begin{equation}
 \mathrm d\tau=\mathrm dt\sqrt{1-v^2/c^2}.
 \label{eq:mech-proper-time}
\end{equation}
匀速钟在两个沿其轨迹的事件间满足 $\Delta t=\gamma_v\Delta\tau$，其中 $\gamma_v=(1-v^2/c^2)^{-1/2}$。比较不同世界线的累计钟时，需要分别积分，不能把某段匀速的关系无条件用于整个往返过程。

设细杆在 $S'$ 中静止且沿运动方向，\keyterm{固有长度（proper length）}为 $\ell_0$。在 $S$ 中测长度必须取两端在 $S$ 中同时的位置，故 $\Delta t=0$，由 $\Delta x'=\Gamma\Delta x$ 得 $\ell=\ell_0/\Gamma$。若选两端不同时刻的事件，它们的坐标差不是这一测量意义下的杆长。

\begin{example}[运动钟与运动杆]
飞船相对实验室速率为 $0.6c$。飞船静止长度为 $100\,\mathrm m$，船上同一只钟两次滴答间隔为 $1\,\mathrm s$。求实验室测得的船长及两次滴答的坐标时间间隔。
\end{example}
\begin{solve}
$\Gamma=1/\sqrt{1-0.6^2}=1.25$。实验室同时测量两端得到长度 $80\,\mathrm m$；两次滴答位于船上同一钟的世界线上，实验室坐标间隔为 $1.25\,\mathrm s$。两个结论选取的是不同事件对，不能把长度收缩和时间膨胀机械地应用于任意两个事件。
\end{solve}

\section{速度变换}
对式~\eqref{eq:mech-lorentz} 求微分，除以 $\mathrm dt'$，得到
\[
 u_x'=\frac{u_x-V}{1-Vu_x/c^2},\qquad
 u_y'=\frac{u_y}{\Gamma(1-Vu_x/c^2)},\qquad
 u_z'=\frac{u_z}{\Gamma(1-Vu_x/c^2)}.
\]
若 $u_x=c,u_y=u_z=0$，变换后仍有 $u_x'=c$。当各速度远小于 $c$ 时，分母接近 $1$，恢复伽利略速度叠加。相对论不是任意小速度下产生完全不同的结果，而是在牛顿近似失效时保留不可忽略的修正。

\section{相对论动量、能量与不变质量}
用固有时定义四速度 $U^\mu=\mathrm d(ct,\mathbf r)/\mathrm d\tau=\gamma_v(c,\mathbf v)$。对于静质量（不变质量）为 $m>0$ 的粒子，相对论力学取四动量 $P^\mu=mU^\mu=(E/c,\mathbf p)$，于是
\begin{equation}
 \mathbf p=\gamma_vm\mathbf v,\qquad E=\gamma_vmc^2,
 \qquad E^2=p^2c^2+m^2c^4.
 \label{eq:mech-relativistic-energy}
\end{equation}
这将动量与能量统一为按洛伦兹变换协同变化的量。$m$ 不随参考系变化，本书不使用“相对论质量”表述。动能为 $K=E-mc^2$，在 $v/c\ll1$ 时展开得 $K=mv^2/2+3mv^4/(8c^2)+\cdots$。无质量粒子没有静止系和沿光线的非零固有时，不能从 $m=0$ 直接代入四速度定义；其能动关系单独取 $E=pc$。

力仍可定义为 $\mathbf F=\mathrm d\mathbf p/\mathrm dt$，但一般不再等于 $m\mathbf a$。对恒定 $m$，微分能动关系并利用 $\mathbf v=\mathbf p c^2/E$，得 $\mathrm dE/\mathrm dt=\mathbf F\cdot\mathbf v$。孤立系统的总四动量守恒，系统的不变质量由总能量和总动量共同确定，可能大于各组成粒子静质量之和，因为内部能量也有贡献。

\begin{example}[对向粒子的完全非弹性碰撞]
两粒子静质量均为 $m$，在质心系中以速率 $v$ 对向运动，碰后结合成静止复合体且无能量逸出。求复合体不变质量 $M$。
\end{example}
\begin{solve}
初总动量为零，末复合体在同一系中静止。能量守恒给出 $Mc^2=2\gamma_vmc^2$，所以 $M=2\gamma_vm>2m$（$v>0$）。初始动能转化为复合体内部能量；若有辐射逸出，必须把辐射的能量和动量加入守恒式。
\end{solve}

\section{习题}
\begin{enumerate}[label=\thechapter.\arabic*,ref=\thechapter.\arabic*]
\item\label{exr:mech-simultaneity} 两事件在 $S$ 中同时，沿 $x$ 方向相距 $d>0$。求在沿正 $x$ 方向运动的 $S'$ 中的时间差，并判断哪个事件先发生。
\item\label{exr:mech-velocity-add} 飞船相对实验室沿正 $x$ 方向以 $0.8c$ 运动，船内粒子相对飞船沿同方向以 $0.6c$ 运动。求实验室测得的粒子速度。
\item\label{exr:mech-photons} 两个光子在某惯性系中能量均为 $\varepsilon$，传播方向夹角为 $\theta$。求两光子系统的不变质量，讨论同向与反向两个极限。
\end{enumerate}

\chapter{拉格朗日与哈密顿方法：选学}
\label{chap:mech-lagrange}

\section{广义坐标与理想约束}
复杂约束下，逐个求约束反力往往掩盖真正需要的运动变量。例如单摆只有一个自由度，用摆角即可描述，不必把二维位置与张力同时作为未知量。\keyterm{广义坐标（generalized coordinates）} $q_1,\ldots,q_n$ 是局部独立地确定构型的参数，不一定有长度量纲。若约束可以写为位置和时间的关系，并局部解成 $\mathbf r_i=\mathbf r_i(q_1,\ldots,q_n,t)$，称为\keyterm{完整约束（holonomic constraint）}。

在固定时刻作与约束相容的无穷小构型变化，得到\keyterm{虚位移（virtual displacement）} $\delta\mathbf r_i=\sum_j(\partial\mathbf r_i/\partial q_j)\delta q_j$。它不是沿真实运动经过 $\mathrm dt$ 后的实际位移。若约束力对所有相容虚位移的总虚功为零，称约束为\keyterm{理想约束（ideal constraint）}；无摩擦固定轨道是典型情形。有些移动约束也可为理想约束，但仍可能通过实际位移做功。

\section{由牛顿方程到拉格朗日方程}
设质点质量恒定，真实力分为主动力与理想约束力。将牛顿方程与虚位移点乘并求和，约束力项消去。定义广义力 $Q_j=\sum_i\mathbf F_i\cdot\partial\mathbf r_i/\partial q_j$，以及动能 $T=\sum_i m_i\dot{\mathbf r}_i^2/2$。由链式法则可验证
\[
 \sum_i m_i\ddot{\mathbf r}_i\cdot\frac{\partial\mathbf r_i}{\partial q_j}
 =\frac{\mathrm d}{\mathrm dt}\frac{\partial T}{\partial\dot q_j}-\frac{\partial T}{\partial q_j}.
\]
其依据是 $\partial\dot{\mathbf r}_i/\partial\dot q_j=\partial\mathbf r_i/\partial q_j$，以及 $\mathrm d(\partial\mathbf r_i/\partial q_j)/\mathrm dt=\partial\dot{\mathbf r}_i/\partial q_j$，后者要求坐标映射足够光滑。各 $\delta q_j$ 独立，故分别得到动力学方程。若主动力来自位置势能 $U(q,t)$，令\keyterm{拉格朗日量（Lagrangian）} $\mathcal L=T-U$，则
\begin{equation}
 \frac{\mathrm d}{\mathrm dt}\frac{\partial\mathcal L}{\partial\dot q_j}
 -\frac{\partial\mathcal L}{\partial q_j}=0,\qquad j=1,\ldots,n.
 \label{eq:mech-lagrange}
\end{equation}
存在非保守广义力时，右端改为 $Q_j^{\mathrm{nc}}$。一般非完整约束不能先任意消去速度约束再照抄这一无约束形式。

\begin{example}[用广义坐标求单摆方程]
对摆长 $\ell$、质量 $m$ 的单摆取广义坐标 $\theta$，从竖直向下方向量起，写出拉格朗日量并导出运动方程。
\end{example}
\begin{solve}
动能 $T=m\ell^2\dot\theta^2/2$，以最低点为势能零点有 $U=mg\ell(1-\cos\theta)$，因此 $\mathcal L=T-U$。代入式~\eqref{eq:mech-lagrange} 得 $m\ell^2\ddot\theta+mg\ell\sin\theta=0$。张力没有出现在方程中，是因为它沿绳方向，而允许虚位移沿圆周切向；求出运动后仍可用径向牛顿方程恢复张力。
\end{solve}

\section{作用量驻值与边界条件}
\keyterm{作用量（action）}为 $S[q]=\int_{t_1}^{t_2}\mathcal L(q,\dot q,t)\,\mathrm dt$。假设轨迹和拉格朗日量足够光滑，比较具有相同端点位置的邻近轨迹，即 $\delta q(t_1)=\delta q(t_2)=0$。变分并分部积分，有
\[
 \delta S=\sum_j\int_{t_1}^{t_2}
 \left(\frac{\partial\mathcal L}{\partial q_j}-\frac{\mathrm d}{\mathrm dt}\frac{\partial\mathcal L}{\partial\dot q_j}\right)\delta q_j\,\mathrm dt.
\]
对所有端点固定的光滑变分，$\delta S=0$ 等价于欧拉--拉格朗日方程。\keyterm{哈密顿原理（Hamilton's principle）}应表述为作用量驻值，驻值轨迹不一定使作用量取得最小值；边界条件改变时，边界项也必须重新处理。

\section{循环坐标与能量函数}
定义\keyterm{广义动量（generalized momentum）} $p_j=\partial\mathcal L/\partial\dot q_j$。若 $\mathcal L$ 不显含某个 $q_j$，则该坐标为\keyterm{循环坐标（cyclic coordinate）}，方程立即给出 $\dot p_j=0$。在有心势的平面运动中，$\mathcal L=m(\dot r^2+r^2\dot\theta^2)/2-U(r)$ 不含 $\theta$，故 $p_\theta=mr^2\dot\theta$ 守恒，正是角动量守恒。

令能量函数 $\mathcal E=\sum_j\dot q_jp_j-\mathcal L$，沿无非保守广义力的解求导，利用拉格朗日方程得到
\begin{equation}
 \frac{\mathrm d\mathcal E}{\mathrm dt}=-\frac{\partial\mathcal L}{\partial t}.
 \label{eq:mech-lagrange-energy}
\end{equation}
因此不显含时间时 $\mathcal E$ 守恒。若坐标变换不显含时间，动能是速度的齐次二次型，且势能不含速度，则 $\sum_j\dot q_j\partial T/\partial\dot q_j=2T$，从而 $\mathcal E=T+U$。移动约束下动能可含一次项和常数项，此时不能无条件把能量函数叫作普通机械能。

\section{哈密顿方程}
若速度 Hessian 矩阵 $(\partial^2\mathcal L/\partial\dot q_i\partial\dot q_j)$ 可逆，便能局部将速度表示为 $\dot q(q,p,t)$。定义\keyterm{哈密顿量（Hamiltonian）} $H(q,p,t)=\sum_jp_j\dot q_j-\mathcal L$。对其全微分，并使用 $p_j=\partial\mathcal L/\partial\dot q_j$ 消去 $\mathrm d\dot q_j$ 项，得到
\begin{equation}
 \dot q_j=\frac{\partial H}{\partial p_j},\qquad
 \dot p_j=-\frac{\partial H}{\partial q_j}.
 \label{eq:mech-hamilton}
\end{equation}
这些一阶方程与正则拉格朗日系统等价。它们把位置与动量作为独立相空间变量，并为后续分析力学、统计力学和量子理论提供共同语言；奇异拉格朗日量需要专门的约束处理，不在本章范围内。

\section{习题}
\begin{enumerate}[label=\thechapter.\arabic*,ref=\thechapter.\arabic*]
\item\label{exr:mech-lagrange-spring} 对一维弹簧振子写出 $\mathcal L,p,H$，并用哈密顿方程恢复牛顿方程。
\item\label{exr:mech-polar-lagrange} 对 $\mathcal L=m(\dot r^2+r^2\dot\theta^2)/2-U(r)$，分别对 $r,\theta$ 写出拉格朗日方程，并与极坐标加速度比较。
\item\label{exr:mech-sliding-pendulum} 质量 $M$ 的小车在水平光滑轨道上运动，车上悬挂质量 $m$、长 $\ell$ 的单摆。取车位移 $x$ 与摆角 $\theta$ 为广义坐标，写出拉格朗日量及守恒的广义动量，并说明其物理意义。
\end{enumerate}

\chapter{计算、实验与模型检验}
\label{chap:mech-computation}

\section{无量纲化与参数的意义}
解析方法揭示结构，数值方法处理难以显式求解的模型，实验则检验模型与实际的联系。三者不能互相替代。以阻尼受迫振子为例，选择特征长度 $x_*>0$，令 $X=x/x_*$、$s=\omega_0t$，方程化为
\[
 \frac{\mathrm d^2X}{\mathrm ds^2}+2\zeta\frac{\mathrm dX}{\mathrm ds}+X=f_*\cos(\nu s),
 \qquad \zeta=\frac\gamma{\omega_0},\quad f_*=\frac{F_0}{m\omega_0^2x_*},\quad \nu=\frac\omega{\omega_0}.
\]
无量纲化把单位选择从物理参数中分离出来，也表明只要 $\zeta,f_*,\nu$ 及无量纲初值相同，不同实验系统在此模型下具有相同的无量纲演化。相似性结论依赖模型本身，不意味着材料非线性或阻力机制也自动相同。

\section{显式欧拉法及其能量误差}
对 $\dot x=v$、$\dot v=a(x,v,t)$，时间步长 $h>0$ 的\keyterm{显式欧拉法（explicit Euler method）}为
\[
 x_{n+1}=x_n+hv_n,\qquad v_{n+1}=v_n+h\,a(x_n,v_n,t_n),\qquad t_n=t_0+nh.
\]
它以旧时刻的导数线性外推。在解足够光滑、右端在所考察有限区间满足适当 Lipschitz 条件且数值轨迹保持有界时，局部单步误差为 $O(h^2)$，固定终止时间的全局误差为 $O(h)$。误差随 $h\to0$ 收敛，并不意味着某个固定步长下的长期行为正确。

对无阻尼谐振子 $a=-\omega_0^2x$，直接展开得
\begin{equation}
 E_{n+1}=(1+h^2\omega_0^2)E_n,
 \qquad E_n=\tfrac12m(v_n^2+\omega_0^2x_n^2).
 \label{eq:mech-euler-energy}
\end{equation}
所以任意固定正步长都使离散能量逐步增长。这种增长是数值方法造成的，不能解释为物理系统出现能量来源。短时曲线看似平滑也不足以证明算法适用。

\section{速度 Verlet 方法}
对恒质量、加速度只依赖位置的系统 $\ddot{\mathbf r}=\mathbf a(\mathbf r)$，可采用\keyterm{速度 Verlet 方法（velocity Verlet method）}。输入为初始位置、速度、步长和终止时间，先计算 $\mathbf a_0=\mathbf a(\mathbf r_0)$，再循环
\begin{align}
 \mathbf r_{n+1}&=\mathbf r_n+h\mathbf v_n+\tfrac12h^2\mathbf a_n,\notag\\
 \mathbf a_{n+1}&=\mathbf a(\mathbf r_{n+1}),\notag\\
 \mathbf v_{n+1}&=\mathbf v_n+\tfrac h2(\mathbf a_n+\mathbf a_{n+1}).
 \label{eq:mech-verlet}
\end{align}
每步输出位置、速度及所需诊断量，到达终止时间或满足接触、碰撞等事件条件时结束或切换模型。对充分光滑且相关导数有界的解，固定时间范围内的全局误差为 $O(h^2)$，每步需一次新的力计算。若一次求力的成本为 $C_F$，$N$ 步的总成本为 $O(NC_F)$；只保留当前状态所需存储与系统自由度成正比。

速度 Verlet 对可分离保守哈密顿系统具有辛结构和时间可逆性，但一般不精确守恒原始机械能。对谐振子，其位置递推为 $x_{n+1}-2x_n+x_{n-1}=-h^2\omega_0^2x_n$，特征根位于单位圆且不重合的稳定范围为 $0<h\omega_0<2$。边界 $h\omega_0=2$ 对一般初值可能出现线性增长，不能包含在稳健的稳定区间内。速度相关阻力、非光滑接触和时变力需要相应修改算法。

\section{步长收敛与守恒量诊断}
计算前应明确输出目标：轨道形状、终点位置、振动周期和长期能量漂移可能对步长有不同要求。对同一物理终止时间，分别用 $h,h/2,h/4$，比较同一时刻的解。若方法阶数为 $p$，且已进入误差主项支配的渐近范围，应有
\[
 \frac{\|z_h-z_{h/2}\|}{\|z_{h/2}-z_{h/4}\|}\approx2^p.
\]
这是一种经验诊断，分母接近零、舍入误差占优、事件时刻未对齐或解不光滑时都可能失效。保守系统可检查相对能量误差和角动量误差；耗散系统则应检查正确的耗散方向及功率平衡，而非要求能量不变。

\begin{example}[引力轨道的计算方案]
对固定力心模型 $\ddot{\mathbf r}=-GM\mathbf r/r^3$，设计一个可核查的数值方案。
\end{example}
\begin{solve}
先给定 $r_0>0$ 和切向初速，并用式~\eqref{eq:mech-effective-potential} 对应的单质点能量及角动量判断预期轨道类型。以 $r_0$ 和 $\sqrt{r_0^3/(GM)}$ 无量纲化，再用速度 Verlet 推进；每步记录 $r$、$E=mv^2/2-GMm/r$ 及 $\mathbf L=m\mathbf r\times\mathbf v$。用圆轨道的解析周期校准，并用步长减半检查终点误差。近心点处变化最快，应使步长足以分辨该段；若轨道趋近 $r=0$ 的奇点或实体表面，停止并重新检查物理模型，不能让算法跨过奇点后继续解释结果。
\end{solve}

\section{以单摆测量为例的实验论证}
实验设计应先明确测量的是哪个模型参数。单摆测 $g$ 时，测量摆长应从悬点到摆球质心，周期可由连续 $N$ 次完整振动的总时间除以 $N$ 估计，以减小人工计时误差占比。选择数个摆长，拟合 $T^2=a\ell+b$；理想小角度模型预测 $a=4\pi^2/g$ 且 $b=0$。显著非零截距可能提示长度零点偏差或模型不足，不能在未分析残差前直接强制截距为零。

回归时还需考虑横轴测量误差、重复测量相关性和幅度变化。大摆幅使周期增加，空气阻力和支点摩擦改变衰减与周期，摆球和绳的有限尺寸改变理想质点模型。应分别报告重复测量离散性、仪器分辨率和模型修正；“测量值接近公认值”并不能单独证明操作与模型均正确。可信的结论应附带适用范围、估计误差及可复查的测量方法。

\section{习题}
\begin{enumerate}[label=\thechapter.\arabic*,ref=\thechapter.\arabic*]
\item\label{exr:mech-euler-proof} 推导式~\eqref{eq:mech-euler-energy}，并说明在固定终止时间 $t_*=Nh$、$h\to0$ 时为什么能量误差仍可趋于零。
\item\label{exr:mech-verlet-frequency} 对速度 Verlet 的谐振子位置递推，设 $x_n=e^{\mathrm i n\vartheta}$，推导离散角频率 $\widetilde\omega=\vartheta/h$ 与真实 $\omega_0$ 的关系，并求其小步长展开至 $h^2$ 项。
\item\label{exr:mech-pendulum-bias} 单摆实验真实摆幅为小量 $\theta_0$，但仍用 $g_{\mathrm{est}}=4\pi^2\ell/T^2$ 估计重力加速度。求相对偏差的首项，判断偏高还是偏低。
\end{enumerate}

\appendix
\chapter{习题解答}
\label{app:mech-solutions}

\section{物理量与数学工具}
\subsection*{习题~\ref{exr:mech-dimension}}
由 $[b]=[F]/[v]$ 得 $[b]=\mathsf M\mathsf T^{-1}$；由 $[c_{\mathrm d}]=[F]/[v]^2$ 得 $[c_{\mathrm d}]=\mathsf M\mathsf L^{-1}$。于是 $[m/b]=\mathsf T$。在线性阻力模型中，$m/b$ 还是速度偏离终端值按指数衰减的特征时间；这个动力学解释需由运动方程得到，量纲本身只能识别它是时间量。

\subsection*{习题~\ref{exr:mech-cross}}
点积为零，叉积为 $(0,0,6)\,\mathrm{N\,m}$。后者是关于原点的力矩。点积虽然有功的量纲，但位置矢量 $\mathbf r$ 并非题目给定的实际位移；没有运动路径，就不能把该点积当成过程中所做的功。恒力功应是 $\mathbf F\cdot\Delta\mathbf r$。

\subsection*{习题~\ref{exr:mech-error}}
偏导数为 $\partial g/\partial\ell=g/\ell$、$\partial g/\partial T=-2g/T$。代入式~\eqref{eq:mech-uncertainty} 得
\[
 \frac{u_g}{g}\approx\sqrt{\left(\frac{u_\ell}{\ell}\right)^2+4\left(\frac{u_T}{T}\right)^2}.
\]
摆幅过大导致小角度周期模型的系统性偏差。重复计时能减小部分随机不确定度，却不能消除此模型偏差。

\section{质点运动学}
\subsection*{习题~\ref{exr:mech-path}}
速度为 $v=t(2A-3Bt)$，非零转向时刻为 $t_m=2A/(3B)$，其位置 $x_m=4A^3/(27B^2)$。区间起点与终点的位置均为零，因此位移为零。运动先前进至 $x_m$ 再返回，路程为 $2x_m=8A^3/(27B^2)$；直接对速度积分只能得到位移，应对速率积分才能得到路程。

\subsection*{习题~\ref{exr:mech-circle}}
角速度为 $\alpha t$，故切向加速度大小 $a_{\mathrm t}=R\alpha$，法向加速度大小 $a_{\mathrm n}=R\alpha^2t^2$。二者正交，所以 $a=R\sqrt{\alpha^2+\alpha^4t^4}$。两者相等要求 $\alpha t^2=1$，正根为 $t=1/\sqrt\alpha$；此处 $\alpha$ 的量纲为时间的负二次方。

\subsection*{习题~\ref{exr:mech-river}}
最短时间要求船相对水的速度全部用于垂直河岸方向，故 $t_{\min}=d/v$，随水向下游漂移 $ud/v$。若要求到达正对岸，船相对水的上游分量应为 $u$，剩余垂岸分量为 $\sqrt{v^2-u^2}$，所以 $t_{\mathrm{opposite}}=d/\sqrt{v^2-u^2}$。相对水的航向与垂岸方向夹角满足 $\sin\theta=u/v$；题设 $v>u$ 保证可在有限时间到达正对岸。

\section{牛顿动力学}
\subsection*{习题~\ref{exr:mech-elevator}}
向上取正，$N-mg=ma$，得 $N=m(g+a)$。保持接触须 $N\ge0$，即 $a\ge-g$。当 $a=-g$ 时支持力为零，处于自由落体的失重极限；若电梯地板向下加速度更大而无其他约束，人会离开地板，不能继续假定相对地板静止。

\subsection*{习题~\ref{exr:mech-pull}}
竖直方向无加速度时 $N=mg-F\sin\alpha$，需要 $F\sin\alpha\le mg$。物体已向右滑动，摩擦向左，于是
\[
 a_x=\frac{F\cos\alpha-\mu_k(mg-F\sin\alpha)}{m}.
\]
若该值为负，物体正在减速；若速度降至零，应重新判断静摩擦状态。若支持力计算为负，接触假设失效，须按离面后的二维运动求解。

\subsection*{习题~\ref{exr:mech-drag-initial}}
解一阶方程得到 $v(t)=v_\infty+(v_0-v_\infty)e^{-t/\tau}$。若 $v_0>v_\infty$，向下速度单调减小并趋于 $v_\infty$。若 $v_0<0$，物体初始向上，速度分量单调增加，经过 $t_{\mathrm{turn}}=\tau\ln[(v_\infty-v_0)/v_\infty]$ 变为零，随后下落并趋于终端速度。阻力 $-bv$ 自动随速度变号，故同一线性方程可跨越转向点使用。

\section{功与能量}
\subsection*{习题~\ref{exr:mech-potential}}
由 $U'(x)=4ax^3-2bx=0$ 得 $x=0$ 与 $x=\pm\sqrt{b/(2a)}$。二阶导数 $U''=12ax^2-2b$，所以原点处为 $-2b<0$，不稳定；另两点处为 $4b>0$，稳定，小振动角频率为 $\sqrt{4b/m}$。这里稳定性判断针对无耗散一维保守模型。

\subsection*{习题~\ref{exr:mech-power}}
由 $\mathrm d(mv^2/2)/\mathrm dt=P$，取正向运动分支得 $v(t)=\sqrt{v_0^2+2Pt/m}$。积分可得
\[
 x(t)-x(0)=\frac{m}{3P}\left[\left(v_0^2+\frac{2Pt}{m}\right)^{3/2}-v_0^3\right].
\]
力为 $F=P/v$；若 $v_0=0$ 且要求从开始瞬间即保持正的恒功率，力在零时刻发散。实际驱动必须在低速阶段受力或转矩限制，之后才可近似恒功率。

\subsection*{习题~\ref{exr:mech-timed-u}}
运动方程为 $m\ddot x=-k(t)x$，于是 $\dot E=m\dot x\ddot x+kx\dot x+\dot kx^2/2=\dot kx^2/2$。即使 $k$ 缓慢改变，该项一般也不为零；缓慢只可能使某些误差或变化率较小，不能推出严格守恒。改变弹簧参数的外部机构可以与振子交换能量。

\section{动量与质点系}
\subsection*{习题~\ref{exr:mech-com}}
积分得 $M=\int_0^\ell\lambda_0(1+x/\ell)\,\mathrm dx=3\lambda_0\ell/2$，一阶质量矩为 $\int_0^\ell x\lambda(x)\,\mathrm dx=5\lambda_0\ell^2/6$。故 $x_C=5\ell/9>\ell/2$，质心向线密度较大的一端偏移。

\subsection*{习题~\ref{exr:mech-restitution}}
由式~\eqref{eq:mech-collision1} 得 $v_1=0$、$v_2=u/2$。末动能为 $\tfrac12(2m)(u/2)^2=mu^2/4$，初动能为 $mu^2/2$，因此损失比例为 $1/2$。总动量末值为 $2m(u/2)=mu$，与初值相同，可作为核查。

\subsection*{习题~\ref{exr:mech-rocket-gravity}}
向上取正，式~\eqref{eq:mech-rocket} 给出 $\dot v=-g+u_eq/(m_0-qt)$，因此
\[
 v(t)=v_0+u_e\ln\frac{m_0}{m_0-qt}-gt.
\]
公式用于喷气持续且质量为正的区间；若干质量为 $m_f>0$，则 $0\le t\le(m_0-m_f)/q$，并且该阶段应足够短、空间变化足够小以支持恒定 $g$ 近似。初始净加速度向上要求 $u_eq>m_0g$。若火箭仍受发射台支撑，在离台前还需计入支持力。

\section{角动量}
\subsection*{习题~\ref{exr:mech-straight-L}}
位置为 $\mathbf r=(x,b,0)$，动量为 $(mv,0,0)$，故 $\mathbf L_O=(0,0,-mbv)$，与 $x$ 无关。角动量不要求轨迹为圆；只要运动直线不经过原点，即 $b\ne0$，角动量可以非零。无外力矩使它保持恒定。

\subsection*{习题~\ref{exr:mech-couple}}
合力为零，两力矩均沿正 $z$ 方向，总力矩为 $2aF\mathbf e_z$。把取矩点平移 $\mathbf d$ 后，总力矩变为 $\sum_i(\mathbf r_i-\mathbf d)\times\mathbf F_i=\boldsymbol\tau_O-\mathbf d\times\sum_i\mathbf F_i=\boldsymbol\tau_O$。因此力偶矩不依赖取矩点。

\subsection*{习题~\ref{exr:mech-stick-impact}}
碰撞时间短，重力角冲量可忽略；轴的未知冲量关于轴的力矩为零。碰前关于轴的角动量大小为 $mu\ell$，碰后转动惯量为 $(M/3+m)\ell^2$，所以
\[
 \omega_+=\frac{mu}{(M/3+m)\ell}.
\]
系统线动量一般不守恒，因为轴有外冲量；碰撞机械能也不守恒，因为弹丸嵌入。适用的是关于固定轴的角动量守恒。

\section{刚体}
\subsection*{习题~\ref{exr:mech-rod-I}}
取距端点 $x$，$\mathrm dm=(M/\ell)\,\mathrm dx$，则 $I_O=(M/\ell)\int_0^\ell x^2\,\mathrm dx=M\ell^2/3$。由平行轴定理，$I_O=M\ell^2/12+M(\ell/2)^2=M\ell^2/3$，两种方法一致。

\subsection*{习题~\ref{exr:mech-disc-roll}}
实心圆柱有 $I_C=MR^2/2$。纯滚动时 $Mgh=Mv_C^2/2+I_Cv_C^2/(2R^2)=3Mv_C^2/4$，故 $v_C=\sqrt{4gh/3}$。若斜面倾角为 $\alpha$，式~\eqref{eq:mech-rolling} 给出所需 $f=Mg\sin\alpha/3$，因此 $\mu_{s,\min}=\tan\alpha/3$。若摩擦不足，前述能量分配中的纯滚动约束不成立。

\subsection*{习题~\ref{exr:mech-physical-pendulum}}
以质心在轴下方的稳定位置为 $\theta=0$，方程为 $I_O\ddot\theta=-Mgd\sin\theta$。小角度线性化得 $T=2\pi\sqrt{I_O/(Mgd)}$。均匀细杆绕端点时 $I_O=M\ell^2/3$、$d=\ell/2$，故 $T=2\pi\sqrt{2\ell/(3g)}$；这相当于摆长为 $2\ell/3$ 的小角度单摆。

\section{万有引力}
\subsection*{习题~\ref{exr:mech-satellite}}
令 $r_2/r_1=4$，在卫星质量均可忽略、中心质量相同的条件下，$v\propto r^{-1/2}$、$T\propto r^{3/2}$，所以 $v_2/v_1=1/2$、$T_2/T_1=8$。若两个伴星质量不可忽略，周期关系须用各自的两体质量和，不能仅比较半径。

\subsection*{习题~\ref{exr:mech-uniform-earth}}
以球心为原点、隧道方向为 $x$ 轴，球内引力沿球心方向且与距离成正比，故 $\ddot x=-(GM/R^3)x$。周期为 $T=2\pi\sqrt{R^3/(GM)}$；从表面静止释放时 $x(t)=R\cos(\sqrt{GM/R^3}\,t)$。这一理想周期等于贴近球面圆轨道的周期，但真实天体密度不均匀时一般不保持这种简谐规律。

\subsection*{习题~\ref{exr:mech-orbit-energy}}
各天体绕质心的半径为 $d/2$，向心方程为 $mv^2/(d/2)=Gm^2/d^2$，所以 $v=\sqrt{Gm/(2d)}$。周期为 $T=\pi d/v=2\pi\sqrt{d^3/(2Gm)}$，符合式~\eqref{eq:mech-kepler-third}。总动能为 $mv^2=Gm^2/(2d)$，势能为 $-Gm^2/d$，故总能量为 $-Gm^2/(2d)$。

\section{非惯性系}
\subsection*{习题~\ref{exr:mech-fictitious}}
惯性系中物体有指向轴线的向心加速度，真实合力必须提供 $m\Omega^2r$。转盘系中物体静止，真实力与离心惯性力平衡。向心力是惯性系中真实合力的径向部分，离心惯性力是换参考系后的方程项，二者作用于同一个研究对象，不是一对作用反作用力。若考察物体对转盘的真实反作用，那是另一个物体所受的接触力。

\subsection*{习题~\ref{exr:mech-turntable}}
支持力为 $mg$，静摩擦需提供 $m\Omega^2r$，因此不滑动条件是 $\Omega^2r\le\mu_sg$。若 $r>0$，等价于 $|\Omega|\le\sqrt{\mu_sg/r}$。这里角速度恒定；启动过程中若有角加速度，还需考虑切向加速度引起的摩擦需求。

\subsection*{习题~\ref{exr:mech-coriolis}}
由右手规则，$\boldsymbol\Omega\times\mathbf v'=\Omega u\mathbf e_y$，故 $\mathbf F_{\mathrm{Cor}}=-2m\Omega u\mathbf e_y$。对任意 $\mathbf v'$，三重积 $\mathbf v'\cdot(\boldsymbol\Omega\times\mathbf v')=0$，所以其功率为零。这一结论允许其改变运动方向，不能推断运动轨迹不受影响。

\section{振动}
\subsection*{习题~\ref{exr:mech-spring-vertical}}
从自然长度位置取向下坐标 $x$，方程为 $m\ddot x=mg-kx$。平衡伸长为 $x_{\mathrm{eq}}=mg/k$，令 $\xi=x-mg/k$ 后得到 $\ddot\xi+(k/m)\xi=0$。初值为 $\xi(0)=-mg/k$、$\dot\xi(0)=0$，故最大伸长为 $2mg/k$，周期为 $2\pi\sqrt{m/k}$。重力改变平衡位置，不改变这个线性模型的固有频率。

\subsection*{习题~\ref{exr:mech-log-decrement}}
相邻同向峰对应相同的振动相位，相隔 $T_d=2\pi/\omega_d$，其包络之比为 $e^{\gamma T_d}=q$。故 $\gamma=\ln q/T_d$、$\omega_d=2\pi/T_d$，以及 $\omega_0=\sqrt{\omega_d^2+\gamma^2}$。取相邻绝对值峰时，时间间隔只有半个振动周期，不能混用两种峰值约定。

\subsection*{习题~\ref{exr:mech-modes}}
令 $\omega_+=\sqrt{k/m}$、$\omega_-=\sqrt{(k+2\kappa)/m}$，初值给出 $q_+(0)=q_-(0)=a$，两者初速度为零，因此
\[
 x_1=\frac a2(\cos\omega_+t+\cos\omega_-t),\qquad
 x_2=\frac a2(\cos\omega_+t-\cos\omega_-t).
\]
若 $\kappa\ll k$，则两频率接近，可用和差化积看出快振荡乘慢包络，形成清晰的拍现象。包络绝对值的重复周期为 $2\pi/|\omega_--\omega_+|$；两频率相距较大时，仍为两模叠加，但不宜用清晰慢包络的近似图像解释。

\section{机械波}
\subsection*{习题~\ref{exr:mech-wave-data}}
由定义，$\Lambda=2\pi/k_w=0.5\,\mathrm m$，$f=\omega/(2\pi)=10\,\mathrm{Hz}$，波速 $c_w=\omega/k_w=5\,\mathrm{m\,s^{-1}}$。质点横向速度为 $y_t=A\omega\sin(k_wx-\omega t)$，最大值为 $0.04\pi\,\mathrm{m\,s^{-1}}$，约 $0.126\,\mathrm{m\,s^{-1}}$，与波速明显不同。

\subsection*{习题~\ref{exr:mech-fixed-free}}
取固定端 $x=0$，空间振型 $X=A\sin kx$。自由端条件 $X'(\ell)=Ak\cos k\ell=0$，非零模要求 $k\ell=(2n-1)\pi/2$，其中 $n=1,2,\ldots$。因此 $f_n=(2n-1)c_w/(4\ell)$，只有基频的奇数倍。$k=0$ 不能产生满足固定端条件的非零自由平移模。

\subsection*{习题~\ref{exr:mech-wave-energy}}
令 $k_n=n\pi/\ell$，则 $\omega_n^2=\mathcal T k_n^2/\lambda$。利用 $\int_0^\ell\sin^2(k_nx)\,\mathrm dx=\int_0^\ell\cos^2(k_nx)\,\mathrm dx=\ell/2$，得到
\[
 K=\frac{\lambda\ell A^2\omega_n^2}{4}\sin^2\omega_nt,\qquad
 U=\frac{\mathcal T\ell A^2k_n^2}{4}\cos^2\omega_nt.
\]
两系数相同，故 $E=\lambda\ell A^2\omega_n^2/4$ 为常量。固定端位移和速度为零，不向外传递功率，与能量守恒一致。

\section{连续介质}
\subsection*{习题~\ref{exr:mech-rods}}
串联时两杆受相同轴向力 $F$，总伸长为 $\Delta\ell=F\ell_1/(Y_1A)+F\ell_2/(Y_2A)$。故 $k_{\mathrm{eq}}=A/(\ell_1/Y_1+\ell_2/Y_2)$。该结果等价于两个线性弹簧的柔度相加，要求接头变形可忽略且两杆都在线弹性范围内。

\subsection*{习题~\ref{exr:mech-dam}}
令 $z$ 为到液面的向下深度，净压强为 $\rho gz$，所以 $F=\int_0^h\rho gzb\,\mathrm dz=\rho gbh^2/2$。关于液面水平轴的力矩为 $\int_0^hz(\rho gzb)\,\mathrm dz=\rho gbh^3/3$，故合力作用深度为 $z_*=2h/3$。它深于面积中心 $h/2$，因为下部压强更大。

\subsection*{习题~\ref{exr:mech-pipe}}
由连续性，$v_1=Q/A_1$、$v_2=Q/A_2$。伯努利方程给出
\[
 p_1-p_2=\frac12\rho Q^2\left(\frac1{A_2^2}-\frac1{A_1^2}\right).
\]
当 $A_2<A_1$，速度增加且 $p_2<p_1$。这是同一流线上定常理想流动的结果；实际管道有黏性损失时还需增加能量损失项。

\section{狭义相对论}
\subsection*{习题~\ref{exr:mech-simultaneity}}
将事件 $2$ 取在事件 $1$ 的正 $x$ 方向，$\Delta x=d$、$\Delta t=0$。由式~\eqref{eq:mech-lorentz}，$\Delta t'=-\Gamma Vd/c^2$。当 $V>0$，$\Delta t'<0$，所以正 $x$ 方向的事件 $2$ 在 $S'$ 中先发生。这两个事件间隔为空间样，改变其时间顺序不意味着能改变因果传递顺序。

\subsection*{习题~\ref{exr:mech-velocity-add}}
使用速度变换的逆式，得 $u=(u'+V)/(1+u'V/c^2)=(0.6+0.8)c/(1+0.48)=35c/37\approx0.946c$。结果小于 $c$；直接相加得到 $1.4c$ 是把伽利略变换用于其不适用的高速范围。

\subsection*{习题~\ref{exr:mech-photons}}
每个光子动量大小为 $\varepsilon/c$，所以总动量平方为 $p_{\mathrm{tot}}^2=2(\varepsilon/c)^2(1+\cos\theta)$，总能量为 $2\varepsilon$。系统质量满足
\[
 M^2c^4=(2\varepsilon)^2-c^2p_{\mathrm{tot}}^2
 =4\varepsilon^2\sin^2(\theta/2),
\]
故 $M=2\varepsilon\sin(\theta/2)/c^2$，其中 $0\le\theta\le\pi$。同向时 $M=0$，反向时 $M=2\varepsilon/c^2$。组成粒子各自无质量，不意味着它们组成的系统不变质量必为零。

\section{分析力学}
\subsection*{习题~\ref{exr:mech-lagrange-spring}}
拉格朗日量为 $\mathcal L=m\dot x^2/2-kx^2/2$，广义动量 $p=m\dot x$，哈密顿量为 $H=p^2/(2m)+kx^2/2$。由 $\dot x=\partial H/\partial p=p/m$ 和 $\dot p=-\partial H/\partial x=-kx$，得到 $m\ddot x=-kx$。该哈密顿量与机械能一致，因为坐标变换不显含时间且势能不含速度。

\subsection*{习题~\ref{exr:mech-polar-lagrange}}
径向方程为 $m\ddot r-mr\dot\theta^2+U'(r)=0$，角向方程为 $\mathrm d(mr^2\dot\theta)/\mathrm dt=0$。当 $r>0$ 时，后者展开为 $m(r\ddot\theta+2\dot r\dot\theta)=0$。它们分别对应式~\eqref{eq:mech-polar-a} 的径向与切向分量，真实力为 $-U'(r)\mathbf e_r$。

\subsection*{习题~\ref{exr:mech-sliding-pendulum}}
取竖直向上坐标，摆球位置为 $(x+\ell\sin\theta,-\ell\cos\theta)$，因此
\[
 \mathcal L=\tfrac12(M+m)\dot x^2+m\ell\dot x\dot\theta\cos\theta
 +\tfrac12m\ell^2\dot\theta^2-mg\ell(1-\cos\theta).
\]
拉格朗日量不含 $x$，故 $p_x=(M+m)\dot x+m\ell\dot\theta\cos\theta$ 守恒。这个量等于小车与摆球的总水平动量，与系统无水平外力一致。若将小车位置错误地视为固定，就会遗漏动能中的交叉项和这一守恒关系。

\section{计算与实验}
\subsection*{习题~\ref{exr:mech-euler-proof}}
显式欧拉更新为 $x_{n+1}=x_n+hv_n$、$v_{n+1}=v_n-h\omega_0^2x_n$。展开 $v_{n+1}^2+\omega_0^2x_{n+1}^2$，交叉项抵消，余项为 $(1+h^2\omega_0^2)(v_n^2+\omega_0^2x_n^2)$。固定 $t_*=Nh$ 时，$E_N/E_0=(1+h^2\omega_0^2)^{t_*/h}$，其对数为 $t_*h\omega_0^2+O(h^3)$，故趋于零，能量比趋于 $1$。固定有限时间的收敛与固定步长的长期失真并不矛盾。

\subsection*{习题~\ref{exr:mech-verlet-frequency}}
代入离散指数解，得 $2\cos\vartheta-2=-h^2\omega_0^2$，所以在 $0<h\omega_0<2$ 的主分支上
\[
 \widetilde\omega=\frac2h\arcsin\frac{h\omega_0}2
 =\omega_0+\frac{h^2\omega_0^3}{24}+O(h^4).
\]
数值频率略大于真实频率，单步相位误差虽小，长期仍会累计。这也说明能量误差有界不能保证相位长期精确。

\subsection*{习题~\ref{exr:mech-pendulum-bias}}
由 $T=T_0[1+\theta_0^2/16+O(\theta_0^4)]$，得
\[
 \frac{g_{\mathrm{est}}}{g}=\left(\frac{T_0}{T}\right)^2
 =1-\frac{\theta_0^2}{8}+O(\theta_0^4).
\]
相对偏差首项为 $-\theta_0^2/8$，因此估计偏低。延长计时、增加振动次数只能减少某些计时误差，不能消除这一振幅修正。

\backmatter
\begin{thebibliography}{9}
\bibitem{mit801}
MIT OpenCourseWare. \textit{8.01SC Classical Mechanics: Syllabus}. Fall 2016.
\url{https://ocw.mit.edu/courses/8-01sc-classical-mechanics-fall-2016/pages/syllabus/}.
课程范围、入门层次与力和守恒律两条主线的参考。检索日期：2026年10月11日。

\bibitem{dourmashkin}
Peter Dourmashkin. \textit{Classical Mechanics: MIT 8.01 Course Notes}. MIT OpenCourseWare，在线教材目录（页面注明2022年更新勘误）.
\url{https://ocw.mit.edu/courses/8-01sc-classical-mechanics-fall-2016/pages/online-textbook/}.
用于主题覆盖与进一步阅读指引；本书未按其章节逐段翻译。检索日期：2026年10月11日。

\bibitem{mit8012}
MIT OpenCourseWare. \textit{8.012 Physics I: Classical Mechanics: Syllabus}. Fall 2008.
\url{https://ocw.mit.edu/courses/8-012-physics-i-classical-mechanics-fall-2008/pages/syllabus/}.
本科进阶课程定位及教材信息来源；页面列示 Daniel Kleppner 与 Robert J. Kolenkow，\textit{An Introduction to Mechanics}，McGraw-Hill，1973。检索日期：2026年10月11日。

\bibitem{tsinghua-chen}
陈曦. 教学. 清华大学物理系个人教学主页.
\url{https://info.phys.tsinghua.edu.cn/chenx/zh-CN/teaching/58ed3ef0a066ee36a1ebcead.html}.
用于了解公开讲义的预期读者与教学组织原则，并非清华全部本科力学课程的统一教学大纲。检索日期：2026年10月11日。

\bibitem{pku-curriculum}
北京大学物理学院. 课程体系：夯实基础，因材施教. 2023年5月19日.
\url{https://110.phy.pku.edu.cn/info/1106/1551.htm}.
本科普通物理与理论物理的层次划分、力学课程范围的参考。检索日期：2026年10月11日。
\end{thebibliography}
\end{document}
